% concatenated from main.tex
\documentclass[12pt, oneside,reqno]{amsart}
\usepackage[utf8]{inputenc}
\usepackage{graphicx} % Required for inserting images
\usepackage{amsmath}
\usepackage{amsthm}
\usepackage{amsfonts}
\usepackage{amssymb}
\usepackage{bbm}
\usepackage{mathtools}
\usepackage{fullpage}
\usepackage{tabularx}
\usepackage{multirow}
\numberwithin{equation}{section}
\usepackage{comment}
\usepackage{enumerate}
\usepackage{appendix}
\usepackage{enumitem}
\usepackage{lipsum}
\usepackage[lmargin=1in,rmargin=1in,tmargin=1in,bmargin=1in]{geometry}
\usepackage[ps,all,arc,rotate]{xy}
\usepackage{graphicx, float, epstopdf}
\usepackage{hyperref, color}
\usepackage{centernot}
\usepackage{fancyhdr}
\usepackage{graphics, setspace}
\usepackage{braket}
\usepackage{tikz}
\usepackage{bbm}


%%%
\newcommand{\skipline}{\vspace{11pt}}   			% might have to be preceeded by \\ in the previous line (clarity)
\allowdisplaybreaks[4]                        %allows breaking multiline environments in amsmath commands
\newcommand{\blankline}{\quad\pagebreak[2]} 
%%%

\newtheorem{theorem}{Theorem}[section]
\newtheorem*{thm}{Theorem}
\newtheorem{lemma}[theorem]{Lemma}
\newtheorem{proposition}[theorem]{Proposition}
\newtheorem{corollary}[theorem]{Corollary}
\newtheorem{remark}[theorem]{Remark}

\newcommand{\psum}{\sideset{}{^*}\sum}
\newcommand{\esum}{\sideset{}{^+}\sum}
\newcommand{\osum}{\sideset{}{^-}\sum}

\makeatletter
\newcommand{\doublewidetilde}[1]{{%
  \mathpalette\double@widetilde{#1}%
}}
\newcommand{\double@widetilde}[2]{%
  \sbox\z@{$\m@th#1\widetilde{#2}$}%
  \ht\z@=.9\ht\z@
  \widetilde{\box\z@}%
}
\makeatother


\title{A central limit theorem for partitions involving generalised divisor functions}
\author[1]{Madhuparna Das}
\address{Department of Mathematics, University of Exeter, Exeter, EX4 4QE, United Kingdom}
\email{md679@exeter.ac.uk}
\author[2]{Nicolas Robles}
\address{RAND Corporation, Engineering and Applied Sciences, 1200 S Hayes St, Arlington, VA, 22202, USA}
\email{nrobles@rand.org}
\date{}

\subjclass[2020]{11P82, 05A17, 60F05, 11M41}
\keywords{Central limit theorem, partition function, Dirichlet series with a shifted coefficient, Euler product, saddle-point method, converse mapping}

\begin{document}

\maketitle

\begin{center}
{\it To Bruce Berndt, on the occasion of his 86th birthday}
\end{center}

\begin{abstract}
We define an $f$-restricted partition $p_f(n,k)$ of fixed length $k$ given by the bivariate generating series
\begin{align*}
Q_f(z,u) 
\coloneqq 1+\sum_{n=1}^{\infty}\sum_{k=1}^{\infty} p_f(n,k) u^kz^n
=\prod_{k=1}^{\infty}(1+uz^k)^{\Delta_f(k)},   
\end{align*}
where $\Delta_f(n)=f(n+1)-f(n)$. In this article, we establish a central limit theorem for the number of summands in such partitions when $f(n)=\sigma_r(n)$ denotes the generalised divisor function, defined as $\sigma_r(n)=\sum_{d|n}d^r$ for integer $r\geq 2$. This can be considered as a generalisation of the work of Lipnik, Madritsch, and Tichy, who previously studied this problem for $f(n)=\lfloor{n}^{\alpha}\rfloor$ with $0<\alpha<1$. A key element of our proof relies on the analytic behaviour of the Dirichlet series
\begin{align*}
\sum_{n=1}^{\infty}\frac{\sigma_r(n+1)}{n^s},    
\end{align*}
for $\mathrm{Re}(s)>1$. We study this problem employing the identity involving the Ramanujan sum. Furthermore, we analyse the Euler product arising from the above Dirichlet series by adopting the argument of Alkan, Ledoan and Zaharescu.
\end{abstract}

\section{Introduction}

An arithmetic function $f:\mathbb{N}\to\mathbb{R}_{>0}$ is called multiplicative if $f(mn)=f(m)f(n)$ whenever $(m,n)=1$. An integer partition $p_f(n)$ of a positive integer $n$ into the function $f$ counts the number of ways to write $n$ as follows
\begin{align}\label{partitiondef}
n=\lfloor{f(n_1)}\rfloor+\lfloor{f(n_2)}\rfloor+\cdots+\lfloor{f(n_r)}\rfloor ,   
\end{align}
where $f(n_j)$ is a real positive valued multiplicative function and $n_1\leq n_2\leq\cdots\leq n_r$. It is a well-known fact that one can write the above expression with the generating series 
\begin{align}\label{generatingseries}
{Q_{f}}(z) 
=1+\sum_{n=1}^{\infty}{{p_{f}}(n){z^n}}  
=\prod_{n\in\mathbb{N}^{*}} (1-z^{\lfloor{f(n)}\rfloor})^{-1} \quad (|z|<1).
\end{align}
\noindent
The asymptotic formula for integer partitions, that is, when $f(n)=n$ for all $n\in\mathbb{N}$ in~\eqref{generatingseries}, was first established by Hardy and Ramanujan~\cite{hardyramanujan}. Over the past century, the study of partition asymptotics has attracted the attention of numerous mathematicians; see, for example,~\cite{taylor,madhu,semiprime,vaughan}. The most general framework in this context was studied by Debruyne and Tenenbaum~\cite{debten}.

\medskip

Let $p_{f}(n)$ and $p_{f}(n,r)$ denote, respectively, the ``restricted" partition of $n$ and and the ``restricted" partition of $n$ with fixed length $r$ as defined in~\eqref{partitiondef}, where $n_1\leq n_2\leq\ldots\leq n_r$. Li and Chen~\cite[Theorem 1.1]{lichen2} showed
an alternative approach to analysing the generating series~\eqref{generatingseries}, which is through the study of
\begin{align}\label{generatingserieswithgapfunction}
Q_f(z)
=\sum_{n=0}^{\infty}p_f(n)z^n
=\prod_{n=1}^{\infty}(1-z^n)^{-\Delta_f(n)}\quad (|z|<1).
\end{align}
where 
\begin{align}\label{gnsetdefinition}
\Delta_f(n)=\#\{m: m\in\mathbb{N}\;\text{and}\;\lfloor{
f(m)}\rfloor=n\}. 
\end{align}

%One of the objectives of this paper is to derive an asymptotic formula for the fixed length partition $p_f(n,r)$ for $r\geq 1$ and $f$ is taken to be the generalized divisor function under certain conditions. Nonetheless, our approach can be extended to a broader class of multiplicative functions that satisfy an identity involving the Ramanujan sum.

%To explore this problem, we need to analyze the partition function $p_f(n)$ and $p_f(n,r)$ for a fixed length $r\geq 1$. We later prove that, 

A related problem was studied by Balasubramanian and Luca~\cite{baluluca}, who investigated the number of solutions to $|\{m: m,n\in\mathbb{N}\quad\text{and}\quad m=\mathcal{F}(n)\}|$, where $\mathcal{F}(n)$ denotes the number of unordered factorisation as the products of integer factors greater than 1, by examining the partitions of the form $n=\lfloor{\sqrt{n_1}}\rfloor+\cdots+\lfloor{\sqrt{n_r}}\rfloor$. Later, this result was thoroughly investigated by Li and Chen~\cite{lichen}.

\medskip

Our main objective of this paper is to establish a central limit theorem for the random variable $\varpi_{f,n}$, which counts the number of summands in a random $f$-restricted partition of $n$ (as defined in Section~\ref{mainresults}) with respect to the multiplicative function $f$. This concept was first introduced by Erd\H{o}s and Lehner~\cite{erdoslehner}, where they studied the distribution of ratio $p(n,r)/p(n)$, with $r=(2\pi^2/3)^{-\tfrac{1}{2}}\sqrt{n}\log n+n\sqrt{n}$ as a function of $n$ and $p(n)$ denoting number of integer partitions of $n$. In this context, it is noteworthy that the weighted moments of integer partitions were studied by Richmond~\cite{richmond1,richmond2}. Let $p_{A}(n,r)$ denotes the partition of $n$ into $r$ summands from a given set $A$. The $k$-th moment is then defined by
\begin{align}\label{weightedmoments}
t^k_A(n)
=\sum_{r=1}^{n}r^kp_A(n,r),
\end{align}
for arbitrary fixed $k$ as $n\to\infty$. Furthermore, Richmond remarked that the study of these weighted moments supports the work of Erd\H{o}s and Lehner.



\medskip

In a recent study, Lipnik, Madritsch, and Tichy~\cite{lipniketal} established a central limit theorem for the partition function $p_f(n)$ with $f(n)=\lfloor{n^{\alpha}}\rfloor$ for $0<\alpha<1$. For our main result, we have considered $f$ to be the generalised divisor function $\sigma_r(n)$ for integers $r\geq 2$.
However, our approach is applicable for a wider class of multiplicative functions, which we explore in detail in later sections.  

The proof presented in~\cite{lipniketal} used a probabilistic approach and Curtiss' theorem to obtain the asymptotics of the moment generating function. Furthermore, the application of converse mapping to study the analytic behaviour of the associated Dirichlet series makes their proof simple and elegant. This also allows us to choose a complex weight function in the bivariate generating series and deal with a challenging associated Dirichlet series, as the remaining components of the proof are comparatively less involved than those encountered in prior works on this problem.

\medskip

Studying the limiting distribution of $p_f(n)$ for $f$ being any multiplicative function has a different flavour compared to previous works. Our analysis relies on additive combinatorics, due to the involvement of the gap properties $(|f(n+1)-f(n)|)_{n}$ of the multiplicative function $f$. While we establish a central limit theorem, deriving an asymptotic formula for~\eqref{weightedmoments} in the setting of $f$-restricted partitions would present additional difficulties, largely due to the behaviour of the gap function $(|f(n+1)-f(n)|)_{n}$. However, it appears plausible that an upper bound could be obtained in future work.

%Although obtaining the asymptotic estimate for $p_f(n)$ is somewhat more straightforward than proving the central limit theorem, we employ the saddle-point method for both results, albeit within different frameworks.

\medskip

Throughout this paper $f\sim g$ means ${f(x)}/{g(x)}\to 1$ as $x\to\infty$. Additionally, $f=o(g)$ and $f=O(g)$ denote $\left|{f(x)}/{g(x)}\right|\to 0$ and $|f(x)|\leq C|g(x)|$ as $x\to\infty$, respectively, for a suitable constant $C>0$. 


\section{The Main Results}\label{mainresults}

Let 
\begin{align}\label{ramanujansum} 
c_m(n) 
=\sum_{\substack{b(\bmod{m})\\(b,m)=1}} e\left(\frac{bn}{m}\right), 
\end{align}
denote the Ramanujan sum, and the $r$-th divisor function is given by\footnote{In Section~\ref{EstimateofOmega1f}, we provide a detailed rationale for selecting this particular function within the context of our analysis.}
\begin{align*} 
\sigma_r(n) 
=\sum_{d|n}d^r. 
\end{align*}



\begin{comment}
\begin{theorem}\label{partitiontheorem}
Let $p_f(n)$ denote the weighted integer partition of $n$, where $f$ be the Ramanujan sums or the $k$-th divisor function. Then, as $n\to\infty$, we have
\begin{align*}
p_f(n)\sim    
\end{align*}
\end{theorem}    
\end{comment}


Define
\begin{align}\label{gfunctiondefintion}
\Delta_f(n)
\coloneqq f(n+1)-f(n).
\end{align}

\noindent
Let $f$ be a positive integer valued multiplicative function. We define a $f$-restricted partition $p_f(n)$ such that the fixed-length partition $p_f(n,k)$ is associated with the bivariate generating series
\begin{align}\label{bivariategeneratingseries} 
Q_f(z,u) 
\coloneqq 1+\sum_{n=1}^{\infty}\sum_{k=1}^{\infty} p_f(n,k) u^kz^n
=\prod_{k=1}^{\infty}(1+uz^k)^{\Delta_f(k)}. 
\end{align}

\medskip

Let $\Pi(n)$ denote the collection of $f$-restricted partitions and let $p_f(n)=|\Pi(n)|$ and $p_f(n,k)=|\Pi(n,k)|$ denote the cardinalities of the sets $\Pi(n)$ and $\Pi(n,k)$, respectively. We define the random variable $\varpi_{f,n}$, which counts the number of summands in a random $f$-restricted partition. The probability distribution of $\varpi_{f,n}$ is given by
\begin{align*}
\mathbb{P}(\varpi_{f,n}=k)
=\frac{p_f(n,k)}{p_f(n)}.
\end{align*}

\noindent
We derive the central limit theorem for the random variable $\varpi_{f,n}$ by examining the bivariate generating function $Q_f(z,u)$, which is defined in~\eqref{bivariategeneratingseries}.

\medskip


In the subsequent lemma, we establish the connection between the generating series $Q_f(z,u)$ and the random variable $\varpi_{f,n}$.

\begin{lemma}\label{generatingfunctionlemma}
Consider the bivariate generating series $Q_f(z,u)$ as defined in~\eqref{bivariategeneratingseries}. The following relation holds
\begin{align*}
Q_f(z,u)
=\prod_{k=1}^{\infty}(1+uz^{k})^{\Delta_f(k)}
=1+\sum_{n=1}^{\infty}p_f(n)\mathbb{E}(u^{\varpi_{f,n}})z^n.
\end{align*}
\end{lemma}

\begin{proof}
The first relation follows directly from~\eqref{bivariategeneratingseries}. Furthermore rearranging the terms in~\eqref{bivariategeneratingseries} yields
\begin{align*}
Q_f(z,u)
=&1+\sum_{n=1}^{\infty}\sum_{k=1}^{\infty}p_f(n,k)u^kz^n\\
=&1+\sum_{n=1}^{\infty}p_f(n)z^n\sum_{k=1}^{\infty}\frac{p_f(n,k)}{p_f(n)}u^k\\
=&1+\sum_{n=1}^{\infty}p_f(n)\mathbb{E}(u^{\varpi_{f,n}})z^n,
\end{align*}
thus concluding the lemma.
\end{proof}

Prior to stating our main result, we note that Lipnik et al.~\cite{lipniketal} considered the function $\Delta_f(n)=\lfloor{(n+1)^{\alpha}}\rfloor-\lfloor{n^{\alpha}}\rfloor$ for any real $0<\alpha<1$. We replace their weight function $\Delta_f(n)$ with $\sigma_r(n+1)-\sigma_r(n)$. Furthermore, an analogous result may be obtained using our method for any positive integer valued multiplicative function that satisfies an identity involving the Ramanujan sum.


\begin{theorem}\label{cltforsigma}
Let $f(n)=\sigma_r(n)$ denote the $r$-th divisor function for fixed integer $r\geq 2$. Furthermore, let $\varpi_{f,n}$ be the random variable counting the number of summands in a random $f$-restricted partition of $n$ as defined in \eqref{bivariategeneratingseries}. Then, as $n\to\infty$, $\varpi_{f,n}$ is asymptotically normally distributed with mean $\mu_{f,n}$ and variance $\nu^2_{f,n}$ i.e.,
\begin{align*}
\mathbb{P}\left(\frac{\varpi_{f,n}-\mu_{f,n}}{\nu_{f,n}}<x\right)=\frac{1}{\sqrt{2\pi}}\int_{-\infty}^{x}e^{-{t^2}/{2}}dt+o(1),    
\end{align*}
uniformly for all $x$ as $n\to\infty$. The mean $\mu_{f,n}$ and the variance $\nu^2_{f,n}$ are given by
\begin{align}\label{meansigma}
\mu_{f,n}
=\sum_{k=1}^{\infty}\frac{\Delta_f(k)}{e^{\eta k}+1}
=O(C_{\mu}(r)n^{\frac{r+1}{r+2}})
\end{align}
and
\begin{align}\label{variancesigma}
\nu^2_{f,n}
=\sum_{k=1}^{\infty}\frac{\Delta_f(k)e^{\eta k}}{(e^{\eta k}+1)^2}-\frac{\left(\sum_{k=1}^{\infty}\frac{\Delta_f(k)ke^{\eta k}}{(e^{\eta k}+1)^2}\right)^2}{\sum_{k=1}^{\infty}\frac{\Delta_f(k)k^2e^{\eta k}}{(e^{\eta k}+1)^2}} 
=O(C_{\sigma}(r)n^{\frac{r+1}{r+2}}),
\end{align}
where the constants $C_{\mu}(r)$ and $C_{\sigma}(r)$ are defined in~\eqref{meanconstant} and~\eqref{varianceconstant}, respectively, and $\eta$ is the implicit solution of
\begin{align*}
n=\sum_{k\geq 1}\frac{k}{e^{\eta k}+1}.   
\end{align*}

\noindent
Additionally, the tail of the distribution is given by
\begin{align*}
\mathbb{P}\left(\frac{\varpi_{f,n}-\mu_{f,n}}{\nu_{f,n}}\geq x\right)\leq 
\begin{cases}
e^{-x^2/2}\left(1+O((\log n)^{-3})\right) & \text{if $0\leq x\leq n^{\frac{r+1}{6(r+2)}}/\log n$},\\
e^{-n^{\frac{r+1}{6(r+2)}}x/2}\left(1+O((\log n)^{-3})\right) & \text{if $x\geq n^{\frac{r+1}{6(r+2)}}/\log n$},
\end{cases}
\end{align*}
and an analogous statement holds for $\mathbb{P}\left(\frac{\varpi_{f,n}-\mu_{f,n}}{\nu_{f,n}}\leq -x\right)$.
\end{theorem}

%Using the relation between the divisor function and the Ramanujan sum we establish the following.
\begin{comment}
\begin{theorem}\label{clt}
Let $\varpi_{f,n}$ be the random variable counting the number of summands in a $f$-restricted partition of $n$ as defined in~\eqref{partitiondef}. Then, as $n\to\infty$, $\varpi_{f,n}$ is asymptotically normally distributed with mean $\mu_{f,n}$ and variance $\nu^2_{f,n}$ i.e.,
\begin{align*}
\mathbb{P}\left(\frac{\varpi_{f,n}-\mu_{f,n}}{\nu_{f,n}}<x\right)=\frac{1}{\sqrt{2\pi}}\int_{-\infty}^{x}e^{-{t^2}/{2}}dt+o(1),    
\end{align*}
uniformly for all $x$ as $n\to\infty$. The mean $\mu_{f,n}$ and the variance $\nu^2_{f,n}$ are given by
\begin{align}\label{mean}
\mu_{f,n}
=\sum_{k=1}^{\infty}\frac{\Delta_f(k)}{e^{\eta k}+1}
=\frac{{2}\sqrt{3}}{\pi\sqrt{\upsilon_2(m)}}n^{\frac{1}{2}} & \text{ if }f=c_m(n) \text{ and } \upsilon_i(m)=\prod_{p|m}\left(1-\frac{1}{p^i}\right) 
\end{align}
and
\begin{align}\label{variance}
\nu^2_{f,n}
=b_{q}n^{\frac{1}{2}} & \text{ if }f=c_m(n) \text{ and } \upsilon_i(m)=\prod_{p|m}\left(1-\frac{1}{p^i}\right),
\end{align}
where the constant 
\begin{align*}
b_m
\end{align*}

Finally, the tail of the distribution is given by
\begin{align*}
\mathbb{P}\left(\frac{\varpi_{f,n}-\mu_{f,n}}{\nu_{f,n}}\geq x\right)\leq     
\end{align*}
and an analogous statement holds for $\mathbb{P}\left(\frac{\varpi_{f,n}-\mu_{f,n}}{\nu_{f,n}}\leq -x\right)$.
\end{theorem}
\end{comment}

\subsection{Initial Setup}\label{initialsetup}

In this section, we provide an outline of the proof of Theorem~\ref{cltforsigma}. We adopt the saddle-point approach from~\cite[\S3]{lipniketal} for the initial construction. However, our analysis necessitates a more intricate argument concerning the distribution of multiplicative functions due to the presence of the Dirichlet series with a shifted coefficient. 
%As noted earlier, the proof of Theorem~\ref{partitiontheorem}, provided in Section~\ref{proofofpartitiontheorem}, is more straightforward, as it involves deriving an asymptotic expression for the single-variable generating series (as given in~\eqref{generatingseries}). In contrast, the analysis of the bi-variate generating series~\eqref{bivariategeneratingseries} is more complex. 

\medskip

We begin by observing that, to establish the central limit theorem for the random variable $\varpi_{f,n}$, it is necessary to show that the normalised moment generating function $\mathbb{E}\left(e^{(\varpi_{f,n}-\mu_{f,n})t/\nu_{f,n}}\right)$ converges to $e^{t^2/2}$ as $t \to \infty$. Finally, we estimate the tail of the distribution using the Chernoff bound.

\medskip

Since $\mu_{f,n}/\nu_{f,n}$ is a constant, we can write
\begin{align*}
\mathbb{E}\left(e^{(\varpi_{f,n}-\mu_{f,n})t/\nu_{f,n}}\right)  
=e^{-\mu_{f,n}/\nu_{f,n}}\mathbb{E}\left(e^{\varpi_{f,n}t/\nu_{f,n}}\right).
\end{align*}

\noindent
To achieve our goal we need to obtain the coefficient of $z^n$ in the bivariate generating series $Q(z,u)$. By Lemma~\ref{generatingfunctionlemma}, the probability generating function is expressed as
\begin{align*}
\mathbb{E}(u^{\varpi_{f,n}})
=\frac{[z^n]Q_f(z,u)}{[z^n]Q_f(z,1)}.
\end{align*}
\noindent
Applying Cauchy integral formula we have
\begin{align*}
\mathcal{Q}_{f,n}(u)
\coloneqq [z^n]Q_f(z,u)
=\frac{1}{2\pi i}\oint_{|z|=e^{-\tau}}\frac{Q_f(z,u)}{z^{n+1}}dz.
\end{align*}
\noindent
A standard transformation gives, for $\tau>0$,
\begin{align}\label{mainintegral}
\mathcal{Q}_{f,n}(u)
=\frac{e^{n\tau}}{2\pi}\int_{-\pi}^{\pi}\exp\left(in\theta+F(\tau+i\theta,u)\right)d\theta,
\end{align}
where
\begin{align}\label{logbivariategeneratingseries}
F(\gamma,u)
\coloneqq \log Q_f(e^{-\gamma},u)
=\log\left(\prod_{k=1}^{\infty}(1+ue^{-\gamma k})^{\Delta_f(k)}\right).
\end{align}

\noindent
To tackle the integral in~\eqref{mainintegral}, we utilise the saddle-point technique, which is often referred to as the method of steepest descent. It can be viewed as a coarse version of the circle method. Accordingly, the primary application of this method is to derive estimates for integrals of the form
\begin{align*}
\int_{-\pi}^{\pi} e^{\Phi(\tau+i\theta)}d\theta,   
\end{align*}
for a suitable function $\Phi$. We will now split the integral into two parts. Choose $\theta_n>0$, then we write
\begin{align}\label{integralsplit}
\int_{-\pi}^{\pi} e^{\Phi(\tau+i\theta)}d\theta
=\int_{|\theta|\leq \theta_n} e^{\Phi(\tau+i\theta)}d\theta
+\int_{\theta_n<|\theta|\leq \pi}e^{\Phi(\tau+i\theta)}d\theta.
\end{align}

\noindent
The first part will compute the contribution near the positive real axis, or, equivalently, around $\theta=0$. This will give us the main term, whereas the contribution from the real line i.e, the estimate arising from $e^{\Phi(\tau+i\theta)-\mathrm{Re}(\Phi(\tau))}$ will give us the error term. Using the Taylor expansion up to the third order, the function $\Phi$ can be expressed as

\begin{align}\label{taylorapproximation}
\Phi(\tau+i\theta)
=\Phi(\tau)+i\theta\Phi'(\tau)-\frac{\theta^2}{2}\Phi''(\tau)+O\left(\sup_{\theta}\left|\theta^3\Phi'''(\tau+i\theta)\right|\right).
\end{align}

\noindent
We now choose $\tau$ such that the first-order derivative $\Phi'(\tau)$ vanishes. Thus, the remaining term of the integral in~\eqref{integralsplit} becomes
\begin{align*}
\int_{|\theta|<\theta_n} e^{\Phi(\tau+i\theta)}d\theta
=e^{\Phi(\tau)}\int_{|\theta|<\theta_n} e^{-\tfrac{\theta^2}{2}\Phi''(\tau)}\left(1+O\left(\sup_{\theta}\left|\theta^3\Phi'''(\tau+i\theta)\right|\right)\right) d\theta.
\end{align*}

\noindent
We need to analyse the integrand with the terms involving $\Phi''(\tau)$ and $\Phi'''(\tau)$ in order to show that the aforementioned transformation and estimates are valid. For a detailed explanation, we refer to~\cite[\S3]{lipniketal}. Recall the Mellin transform of a function $h(t)$, defined by
\begin{align*}
h^{*}(s)
\coloneqq\mathcal{M}[h;s]
=\int_{0}^{\infty}h(t)t^{s-1}dt,
\end{align*}
which serves as one of the key ingredients in our analysis. Of particular importance is the rescaling rule (see, for instance,~\cite[Theorem 1]{flajoletetal}), given by
\begin{align*}
\mathcal{M}\left[\sum_{k}\lambda_kh(\mu_kx);s\right]
=h^{*}(s)\sum_{k}\frac{\lambda_k}{\mu_k^s}.
\end{align*}

\noindent
Let $\delta>0$ and set $\delta\leq u\leq \delta^{-1}$. By ``uniformly in $u$", we mean ``uniformly as $\delta\leq u\leq \delta^{-1}$" throughout this paper. Thus, the Mellin transform for the polynomial with respect to $\gamma$ yields
\begin{align*}
\mathcal{M}\left[\sum_{k}\Delta_f(k)\log(1+ue^{-\gamma k});s\right]   
=\mathcal{D}_f(\Delta;s)\mathcal{I}(s,u),
\end{align*}
where
\begin{align}\label{dirichletserieswithG}
\mathcal{D}_f(\Delta;s)
\coloneqq \sum_{n=1}^{\infty}\frac{\Delta_f(n)}{n^s}, 
\end{align}
and
\begin{align*}
\mathcal{I}(s)
=\int_{0}^{\infty} (1+ue^{-\gamma})\gamma^{s-1}d\gamma,      
\end{align*}
represents the Mellin transform of $\gamma\to(1+ue^{-\gamma})$.

\medskip

The principal advantage of the Mellin transform lies not in the transform itself, but in the inverse mapping, which allows us to examine the singularities of the transformed function to obtain the asymptotic expansion\footnotetext{We adopt the term `converse mapping' in Theorem~\ref{cmtheorem} from~\cite{lipniketal}, where it is used in the context of the inverse Mellin transform.}. This step is essential in our case, as we will be working with a Dirichlet series with a shifted multiplicative coefficient.

\begin{theorem}[Converse Mapping~\cite{flajoletetal}]\label{cmtheorem}
Let $f(x)$ be a continuous function on the interval $(0,+\infty)$, and suppose that its Mellin transform $f^{*}(s)$ has a non-empty fundamental strip $\langle\alpha,\beta\rangle$. Assume that $f^{*}(s)$ admits a meromorphic continuation to the larger strip $\langle\gamma,\beta\rangle$, for some $\gamma<\alpha$, with only finitely many poles in that region, and is analytic on the line $\mathrm{Re}(s)=\gamma$. Suppose further that there exists a real number $\eta\in(\alpha,\beta)$ such that
\begin{align}\label{cmcondition1}
f^{*}(s)=O(|s|^{-r})
\end{align}
as $|s|\to\infty$ in $\gamma\leq\mathrm{Re}(s)\leq \eta$, for some $r>1$. If, in addition, $f^*(s)$ admits a singular expansion of the form
\begin{align}\label{cmcondition2}
f^{*}(s)\asymp\sum_{(\xi,k)\in A}\frac{d_{\xi,k}}{(s-\xi)^k} \end{align}
for $s\in\langle\gamma,\alpha\rangle$, then the function admits the following asymptotic expansion as $x\to 0^{+}$:
\begin{align*}
f(x)
=\sum_{(\xi,k)\in A} d_{\xi,k}\left(\frac{(-1)^{k-1}}{(k-1)!}x^{-\xi}(\log x)^{k-1}\right)+O(x^{-\gamma}).
\end{align*}
\end{theorem}

\noindent
Therefore, in our case, we must consider the singularities of the Dirichlet series $\mathcal{D}_f(\Delta;s)$ and $\mathcal{I}(s,u)$. Thus, we are left with the expression

\begin{align*}
\mathcal{I}(s,u)
=\int_{0}^{\infty} (1+ue^{-\gamma})\gamma^{s-1}d\gamma
=\int_{0}^{\infty}\sum_{j=1}^{\infty} \frac{(-u)^{j}}{j}e^{-j\gamma}\gamma^{s-1}d\gamma
=\mathrm{Li}_{s+1}(-u)\Gamma(s),
\end{align*}
where $\mathrm{Li}_{s}(u)$ denotes the polylogarithm, and $\Gamma(s)$ represents the Gamma function. Now, we verify that conditions~\eqref{cmcondition1} and~\eqref{cmcondition2} hold for $\mathrm{Li}_{s}(u)$ and $\Gamma(s)$. An application of Stirling formula yields (see for instance~\cite[\S5 (A.4)]{iwanieckowlaski})
\begin{align*}
|\Gamma(\sigma+it)|=\sqrt{2\pi}|t|^{\sigma-\tfrac{1}{2}}e^{-\pi|t|/2}(1+O_{a,b}(1/t)),   
\end{align*}
where $a\leq\sigma\leq b$ and $|t|\geq 1$. For the polylogarithm function, we adopt the approach outlined in~\cite{lipniketal}. Let $w=-\log z$, and define
\begin{align*}
\Lambda(w)\coloneqq\mathrm{Li}_{r}(e^{-w})
=\sum_{n=1}^{\infty}\frac{e^{-nw}}{n^{r}}.
\end{align*}

\noindent
As this expression represents a harmonic sum, we apply Mellin transform and obtain
\begin{align*}
\Lambda^{*}(s) 
=\int_{0}^{\infty}\sum_{n=1}^{\infty}\frac{e^{-nw}}{n^r}w^{s-1}dw
=\sum_{n=1}^{\infty}\frac{1}{n^{s+r}}\int_{0}^{\infty}e^{-v}v^{s-1}dv
=\zeta(s+r)\Gamma(s),
\end{align*}
for $\mathrm{Re}(s)>\max(0,1-r)$. The Gamma function has simple poles at the negative integers, while $\zeta(s+r)$ has a simple pole at $s=1-r$. Therefore, by applying Theorem~\ref{cmtheorem}, we obtain
\begin{align*}
\mathrm{Li}_{r}(z)
=\Gamma(1-r)w^{r-1}+\sum_{j=0}^{\infty}\frac{(-1)^{j}}{j!}\zeta(r-j)w^{j},
\end{align*}
where 
\begin{align*}
w=\sum_{\ell=1}^{\infty}\frac{(1-z)^{\ell}}{\ell}.    
\end{align*}
\noindent
By applying these estimates in the converse mapping, we derive an asymptotic expression for $\mathcal{Q}_{f,n}$, given by
\begin{align}\label{saddlepointtheorem}
\mathcal{Q}_{f,n}(u)
=\frac{e^{n\tau+F(\tau,u)}}{\sqrt{2\pi B}}(1+O(\tau^{2r/7-1})).
\end{align}
\noindent
Recall the moment-generating function
\begin{align*}
M_n(\theta)=\mathbb{E}(e^{(\varpi_{f,n}-\mu_{f,n})\theta/\nu_{f,n}})
=\exp\left(-\frac{\mu_{f,n}\theta}{\nu_{f,n}}\right)\mathbb{E}(e^{\theta\varpi_{f,n}/\nu_{f,n}})
=\exp\left(-\frac{\mu_{f,n}\theta}{\nu_{f,n}}\right)\frac{\mathcal{Q}_{n}(e^{1/\nu_{f,n}})}{\mathcal{Q}_{n}(1)}.
\end{align*}

\noindent
Next, by applying the Taylor expansion of $M_n(\theta)$, we obtain
\begin{align*}
M_n(\theta)
=\exp\left(\frac{\theta^2}{2}+O(n^{-2r/7(r+2)+1/(r+2)})\right),
\end{align*}
thereby establishing the central limit theorem for $\varpi_{f,n}$.

\section{The Dirichlet series}

In this section, we examine the asymptotic behaviour of the Dirichlet series twisted by the shifted divisor coefficient, which plays a crucial role in deriving the main term. By substituting the expression for $\Delta_f(n)$ into the Dirichlet series given in~\eqref{dirichletserieswithG}, we arrive at
\begin{align}\label{shiteddirichletseries}
\mathcal{D}_f(\Delta;s)
\coloneqq\sum_{n=1}^{\infty}\frac{f(n+1)-f(n)}{n^s}
%=&\sum_{n=1}^{\infty}\frac{f(n+1)}{n^{s}}
%-\sum_{n=1}^{\infty}\frac{f(n)}{n^{s}}\nonumber\\
\eqqcolon A_f(s)-D_f(s),
\end{align}
where 
\begin{align}\label{shifteddirichletseries}
A_f(s)
=\sum_{n=1}^{\infty}\frac{f(n+1)}{n^s}
\end{align}
is the Dirichlet series with a shifted coefficient and 
\begin{align}\label{dirichletseriesdefinition}
D_{f}(s)
=\sum_{n=1}^{\infty}\frac{f(n)}{n^s},
\end{align}
for $\mathrm{Re}(s)>\alpha$ and $\alpha\in\mathbb{R}_{>0}$.

\medskip

Note that for a class of multiplicative functions $f:\mathbb{N}\to\mathbb{N}$ with period 1, i.e. $f(n+1)=f(n)$ for all values of $n$, the analytic properties of $A_f(s)$ are closely analogous to those of $D_f(s)$. One notable example in this context is the Dirichlet character $\chi(n)$ with period 1.  However, determining the analytic behaviour of $A_f(s)$ (such as its poles and residues) directly from the Dirichlet series $D_f(s)$ for general multiplicative functions $f$ remains a difficult problem.
%and a result by Heath-Brown, who proved that
%\begin{align*}
%\sigma_0(n)=\sigma_0(n+1),    
%\end{align*}
%holds for infinitely many values of $n$, where $\sigma_r(n)=\sum_{d|n}d^{r}$ denotes the $r$-th divisor function.
\medskip

The study of the distribution and growth of the gap function $(|f(n+1)-f(n)|)_{n}$ presents an intriguing area of research, leading to significant insights into the distribution of multiplicative functions at primes. In this context, notable contributions have been made by O. Klurman~\cite{klurman} and A. Mangerel~\cite{sacha}. Using the properties of correlations of ``pretentious" functions, Klurman~\cite{klurman} counts the number of ways to write $n=a+b$, where $a,b$ belongs to some multiplicative subsets of $\mathbb{N}$. Recently, Mangerel~\cite{sacha} has explored this type of result in further detail. However, the main result of this paper takes a different approach, connecting instead with counting problems of the type described in~\eqref{gnsetdefinition}.

\medskip

In particular, the Ramanujan sum $c_m(n)$ possesses arithmetic identities with a broad class of multiplicative functions. Furthermore, it can be interpreted as a character, providing significant leverage to deduce the properties of $A_f(s)$ from $D_f(s)$ for those multiplicative functions $f:\mathbb{N}\to\mathbb{N}$ that exhibit an identity with the Ramanujan sum.

For the case of generalised divisor function $f(n)=\sigma_r(n)$, we use the identity (see for instance~\cite{kratzel}) for fixed $r\geq 1$,
\begin{align}\label{kthdivisoridentity}
\sigma_r(n)
=\zeta(r+1)n^r\sum_{m=1}^{\infty}\frac{c_m(n)}{m^{r+1}},
\end{align} 
with the Ramanujan sum given in~\eqref{ramanujansum}.

Note that the case $0\leq r<1$ poses different challenges within the framework of our analysis. Observe that when $f(n)=n^{\beta}$ for $\beta\in\mathbb{R}$, it simplifies to the result in~\cite[Theorem 2]{lipniketal}.

\medskip

The function $f$, considered in this paper satisfies an identity involving the Ramanujan sum $c_m(n)$.  These relations are crucial for handling the Dirichlet series $A_f(s)$ with shifted a coefficient. Moreover, our proof can be extended to encompass any multiplicative functions that exhibit similar identities with the Ramanujan sum. Notably, this method also applies to the von Mangoldt function $\Lambda(n)$. Although $\Lambda(n)$ is not a multiplicative function, it satisfies a similar identity involving the Ramanujan sum. Consequently, the arguments presented in this article are sufficient to establish a central limit theorem, analogous to Theorem~\ref{cltforsigma}, for prime-detecting function that has an identity with $c_m(n)$.

\begin{remark}
In this context, it is important to observe that the Ramanujan sum $c_m(n)$ is a sign-changing function, which implies that integer partitions involving the Ramanujan sum cannot be expressed in terms of~\eqref{partitiondef}. Nevertheless, an analogue of Theorem~\ref{cltforsigma} can be established for the Ramanujan sum, subject with certain restrictions on $m$.

%Nevertheless, the asymptotic behavior and central limit theorem for $c_m(n)$ can be interpreted as a generalization of some previous studies, such as those in~\cite{taylor2}. Consequently, all these studies can be viewed as special cases of our Theorem~\ref{partitiontheorem} and Theorem~\ref{clt} when $f$ is taken to be $c_m(n)$.      
\end{remark}

For Dirichlet character $\chi$ modulo $m$, the generalised Gauss sum is defined as
\begin{align}\label{generalizedgausssum}
c'_{\chi}(n) 
\coloneqq \sum_{\substack{b(\bmod{m})\\(b,m)=1}}\chi(b)e\left(\frac{bn}{m}\right).
\end{align}
Before analysing the Dirichlet series $A_{f}(s)$, we establish some prerequisites.

\begin{lemma}\label{shiftedramanujansumformula}
Let $c_m(n)$ denote the Ramanujan sum for $m\geq 1$ and $n\geq 1$. Then the following identity holds
\begin{align*}
c_m(n+1)
=\frac{\mu(m)}{\varphi(m)}c_m(n)+\frac{1}{\varphi(m)}\sum_{\substack{\chi(\bmod{m})\\ \chi\ne\chi_0}}\tau(\chi)c'_{\bar{\chi}}(n),
\end{align*}
where $\tau(\chi)=c'_{\chi}(1)$ has defined in~\eqref{generalizedgausssum}. %Here, $\mu(m)$ denotes M\"{o}bius function, and $\varphi(m)$ is Euler’s totient function.
\end{lemma}


\begin{proof} 
Considering the Fourier expansion of the additive character (see, for instance,~\cite[\S3.4]{iwanieckowlaski}), we write
\begin{align}\label{fourierexpansion} 
e\left(\frac{b}{m}\right)
=\frac{1}{\varphi(m)}\sum_{\substack{\chi(\bmod{m})}}\Bar{\chi}(b)\tau(\chi), \quad \text{if $(b, m) = 1$}, 
\end{align}
where $\tau(\chi)=\sum_{b(\bmod{m})}\chi(b)e(b/m)$ denotes the Gauss sum associated with the Dirichlet character $\chi$. We simplify the shifted Ramanujan sum as follows

\begin{align}\label{generalizedramanujansum}
c_m(n+1)
=\sum_{\substack{b\bmod{m}\\(b,m)=1}}e\left(\frac{b(n+1)}{m}\right)
=&\sum_{\substack{b\bmod{m}\\(b,m)=1}}e\left(\frac{bn}{m}\right)e\left(\frac{b}{m}\right)\nonumber\\
=&\sum_{\substack{b\bmod{m}\\(b,m)=1}}e\left(\frac{bn}{m}\right)\left(\frac{1}{\varphi(m)}\sum_{\chi(\bmod{m})}\Bar{\chi}(b)\tau(\chi)\right)\nonumber\\
=&\frac{1}{\varphi(m)}\sum_{\chi(\bmod{m})}\tau(\chi)\sum_{\substack{b\bmod{m}\\(b,m)=1}}\Bar{\chi}(b) e\left(\frac{bn}{m}\right),
%=&\frac{1}{\varphi(m)}\sum_{\chi(\bmod{m})}\tau(\chi)\tau(\Bar{\chi})\chi(n),\quad\text{if $(m,n)=1$}\nonumber\\
%=&\frac{m}{\varphi(m)}\sum_{\substack{\chi(\bmod{m})}}{\chi}(-1)\chi(n),\quad\text{if $(m,n)=1$},
%=&m\sum_{d|(n-1,m)}\varphi(d)\mu\left(\frac{m}{d}\right)
\end{align}
where we substituted~\eqref{fourierexpansion} into the second line. We shall isolate the terms associated with principal characters as these are considerably simpler. Consequently, from the preceding expression, we derive
\begin{align*}
c_m(n+1)
=&\frac{1}{\varphi(m)}\left(\sum_{\substack{\chi(\bmod{m})\\ \chi=\chi_0}}\tau(\chi)\sum_{\substack{b\bmod{m}\\(b,m)=1}}\Bar{\chi}(b) e\left(\frac{bn}{m}\right)+\sum_{\substack{\chi(\bmod{m})\\ \chi\ne\chi_0}}\tau(\chi)\sum_{\substack{b\bmod{m}\\(b,m)=1}}\Bar{\chi}(b) e\left(\frac{bn}{m}\right)\right)\nonumber\\
%=&\frac{1}{\varphi(m)}\left(\sum_{\substack{\chi=\chi_0\\\chi(\bmod{m})}}\tau(\chi)\sum_{\substack{b\bmod{m}\\(b,m)=1}}e\left(\frac{bn}{m}\right)+O\left(\sum_{\substack{\chi\ne\chi_0\\\chi(\bmod{m})}}\tau(\chi)c_{\bar{\chi}}(n)\right)\right)\\
=&\frac{\mu(m)}{\varphi(m)}c_m(n)+\frac{1}{\varphi(m)}\sum_{\substack{\chi(\bmod{m})\\ \chi\ne\chi_0}}\tau(\chi)c'_{\bar{\chi}}(n),
\end{align*}
thereby concluding the lemma.
\end{proof}

In the previous lemma, we transformed the shifted Ramanujan sum into a multiplicative structure by using the character sum. This approach facilitates the analysis of the Dirichlet series that arises in the preceding argument.

\begin{lemma}\label{dsklemma}
For a fixed integer $r\geq 1$, consider the double Dirichlet series
\begin{align}\label{dskdefinition}
D_1(s,r)
\coloneqq &\sum_{m=1}^{\infty}\frac{\mu(m)}{\varphi(m)m^{r+1}}\sum_{n=1}^{\infty}\frac{c_m(n)}{n^s},
\end{align}
where $c_m(n)$ denotes the Ramanujan sum. The following identity then holds for $\mathrm{Re}(s)>1$:
\begin{align*}
D_1(s,r)
=\frac{\zeta(s)}{\zeta(s+r+1)}K_r(s),
\end{align*}
where
\begin{align}\label{arsdefinition}
K_r(s)
=C(r)\prod_{p}\left(1+A_{p,r}(s)\right)
\end{align}
with 
\begin{align*}%\label{arpsdefinition}
A_{p,r}(s)
=-\frac{\frac{1}{p^{r+1}(p-1)}}{1+\frac{1}{p^{r+1}(p-1)}}\left(\frac{1}{p^s}-\frac{1}{p^{s+r+1}}+\frac{1}{p^{2s+r+1}}-\frac{1}{p^{2s+2r+2}}+\cdots-\frac{1}{p^{(k+1)(s+r+1)}}+\cdots\right)
\end{align*}
and
\begin{align}\label{crdefinition}
C(r)
\coloneqq\prod_{p}\left(1+\frac{1}{p^{r+1}(p-1)}\right).    
\end{align}
\end{lemma}    


\begin{proof}
Using relation (3.3) of~\cite[\S 3.2]{iwanieckowlaski}, the following expression simplifies as
\begin{align*}
\frac{c_m(n)}{\varphi(m)}
=\frac{1}{\varphi(m)}\frac{\varphi(m)\mu\left({m}/{(m,n)}\right)}{\varphi\left({m}/{(m,n)}\right)}
=\frac{\mu\left({m}/{(m,n)}\right)}{\varphi\left({m}/{(m,n)}\right)}.
\end{align*}

\noindent
By rearranging the terms in~\eqref{dskdefinition} and substituting the above relation, we obtain
\begin{align}\label{dskmodified}
D_1(s,r)
=\sum_{m=1}^{\infty}\frac{\mu(m)}{m^{r+1}}\sum_{n=1}^{\infty}\frac{\mu({m}/{(m,n)})}{\varphi\left({m}/{(m,n)}\right)n^{s}}.
%=\zeta(\mathfrak{s})\sum_{m=1}^{\infty}\frac{\mu(m)(-1)^{\omega(m)}}{m^{r+1+\mathfrak{s}}}\prod_{p|m}\left(\frac{p^{\mathfrak{s}}-1}{p-1}-1\right),
\end{align}

\noindent
Let $\delta=(m,n)$, then we may write $m=k\delta$ and $n=\ell\delta$. Consequently, upon making the change of variables in \eqref{dskmodified}, we obtain
\begin{align*}
D_1(s,r) 
=&\sum_{\delta\geq 1}\mathop{\sum\sum}_{(m,n)=\delta}\frac{\mu(m)}{m^{r+1}}\frac{1}{n^{s}}\frac{\mu({m}/{(m,n)})}{\varphi\left({m}/{(m,n)}\right)}\\
=&\sum_{\delta\geq 1}\mathop{\sum\sum}_{(k,\ell)=1}\frac{\mu(k\delta)}{(k\delta)^{r+1}}\frac{1}{(\ell\delta)^{s}}\frac{\mu(k)}{\varphi(k)}.
\end{align*}
\noindent
Rearranging the order of summation yields
\begin{align*}
D_1(s,r)
=\sum_{k=1}^{\infty}\frac{\mu^2(k)}{\varphi(k)k^{r+1}}\mathop{\sum\sum}_{(\delta,k)=1}\frac{\mu(\delta)}{\delta^{s+r+1}}\sum_{(\ell,k)=1}\frac{1}{\ell^{s}}.
\end{align*}
\noindent
The sum over $\ell$ equals to
\begin{align*}
\zeta(s)\prod_{p|k}\left(1-\frac{1}{p^{s}}\right),
\end{align*}
and the sum over $\delta$ equals to
\begin{align*}
\frac{1}{\zeta(s+r+1)}\prod_{p|k}\left(1-\frac{1}{p^{s+r+1}}\right)^{-1}.
\end{align*}

\noindent
By rearranging the sum over $k$ and invoking multiplicativity, we deduce that
\begin{align}\label{dskeulerproductform}
D_1(s,r)
=\frac{\zeta(s)}{\zeta(s+r+1)}\prod_p\left(1+\frac{1}{p^{r+1}(p-1)}\frac{1-p^{-{s}}}{1-p^{-(s+r+1)}}\right)
=\frac{\zeta(s)}{\zeta(s+r+1)}K_r(s).
\end{align}

\noindent
Observe that as $\mathrm{Re}(s)\to\infty$, the factors $1-p^{-s}$ and $1-p^{-(s+r+1)}$ tend to 1 for any fixed integer $r\geq 1$. In contrast, the factor 
${1}/{p^{r+1}(p-1)}$ remains unchanged, as it is independent of $s$. Hence, we normalise the Euler product $K_r(s)$ as
\begin{align*}
K_r(s)
=&\prod_p\frac{\left(1+\frac{1}{p^{r+1}(p-1)}\frac{1-p^{-{s}}}{1-p^{-(s+r+1)}}\right)}{1+\frac{1}{p^{r+1}(p-1)}}\left(1+\frac{1}{p^{r+1}(p-1)}\right)\\
=&{C(r)}\prod_{p}\frac{\left(1+\frac{1}{p^{r+1}(p-1)}(1+B_{p,r}(s))\right)}{1+\frac{1}{p^{r+1}(p-1)}},
\end{align*}
where $C(r)$ is given in~\eqref{crdefinition}, and
\begin{align}\label{bprs}
B_{p,r}(s)
=\frac{1-\frac{1}{p^s}}{1-\frac{1}{p^{s+r+1}}}-1
=&\left(1-\frac{1}{p^s}\right)\left(1+\frac{1}{p^{s+r+1}}+\frac{1}{p^{2(s+r+1)}}+\cdots\right)-1\nonumber\\
=&-\left(\frac{1}{p^s}-\frac{1}{p^{s+r+1}}+\frac{1}{p^{2s+r+1}}-\cdots-\frac{1}{p^{(k+1)(s+r+1)}}+\cdots\right).
\end{align}

\noindent
In this context, it is worth noting that the decomposition argument employed herein was adopted from \cite[\S3]{alexandru}. Thus, the Euler product in~\eqref{dskeulerproductform} can be expressed as
\begin{align*}
K_r(s)=C(r)\prod_{p}(1+A_{p,r}(s)) ,   
\end{align*}
where
\begin{align*}
A_{p,r}(s)
=&\frac{1+\left(\frac{1}{p^{r+1}(p-1)}\left(1+B_{p,r}(s)\right)\right)}{1+\frac{1}{p^{r+1}(p-1)}}-1\\
=&-\frac{\frac{1}{p^{r+1}(p-1)}}{1+\frac{1}{p^{r+1}(p-1)}}\left(\frac{1}{p^s}-\frac{1}{p^{s+r+1}}+\frac{1}{p^{2s+r+1}}-\frac{1}{p^{2s+2r+2}}+\cdots-\frac{1}{p^{(k+1)s+r+1}}+\cdots\right).
\end{align*}
By~\eqref{bprs}, the above assertion follows, thereby concluding the lemma.
\end{proof}

\begin{remark}
Note that the constant $C(r)$ defined in~\eqref{crdefinition} may be regarded as the generalised totient summatory constant. In particular, when $r=1$, we obtain
\begin{align*}
C(1)=\prod_{p}\left(1+\frac{1}{p^2(p-1)}\right)=1.339784\ldots,   
\end{align*}
a result that can be deduced from the well-known Landau totient constant
\begin{align*}
A
=\zeta(2)C(1)
=\zeta(2)\prod_{p}\left(1+\frac{1}{p^2(p-1)}\right)
=2.20386\ldots.  
\end{align*}
\end{remark}

Let $s=\sigma+it$ with $\sigma, t\in\mathbb{R}$. We have
\begin{align*}
\left|1-\frac{1}{p^{s+r+1}}\right|  
\geq 1-\frac{1}{2^{\sigma+r+1}},\quad\mathrm{and}\quad
\left|\frac{1}{p^s}\right|
= \frac{1}{p^{\sigma}},
\end{align*}
then
\begin{align*}
|B_{p,r}|
=&\left|\frac{1}{p^s}-\frac{1}{p^{s+r+1}}+\frac{1}{p^{2s+r+1}}-\frac{1}{p^{2s+2r+2}}+\cdots-\frac{1}{p^{(k+1)s+r+1}}+\cdots\right|
<&\frac{\frac{1}{p^{\sigma}}}{1-\frac{1}{p^{\sigma+r+1}}}.
\end{align*}
\noindent
Additionally, for a fixed integer $r\geq 1$,
\begin{align*}
{\frac{1}{p^{r+1}(p-1)}}\left({1+\frac{1}{p^{r+1}(p-1)}}\right)^{-1}<\frac{1}{p^{r+1}}.    
\end{align*}
\noindent
so that
\begin{align*}
\sum_{p}|A_{p,r}(s)|
<\frac{2^{\sigma+r+1}}{2^{\sigma+r+1}-1}\sum_{p}\frac{1}{p^{\sigma+r+1}}
<\infty.
\end{align*}
\noindent
It follows that the Euler product $K_r(s)$ converges absolutely (see, for instance,~\cite[\S1.42]{titchmarsh}) and defines an analytic function on the half-plane $\mathrm{Re}(s)>1$. Furthermore, for any fixed $\sigma_0>-r$,
\begin{align*}
|K_r(s)|
\leq C(r)\prod_{p}\left(1+|A_{p,r}(s)|\right)
\leq C(r)\prod_{p}\left(1+\frac{2^{\sigma_0+r+1}}{(2^{\sigma_0+r+1}-1)p^{\sigma_0+r+1}}\right).
\end{align*}

\noindent
Hence, $K_r(s)$ is uniformly bounded on the half-plane $\mathrm{Re}(s)>\sigma_0$. Since it is analytic for $\mathrm{Re}(s)>1$, the Dirichlet series $D_1(s,r)$ admits a meromorphic continuation in this region. Moreover, because $\zeta(s)$, it has a simple pole at point $s=1$ and $1/\zeta(s+r+1)$ has no zeros on the half-plane where $\mathrm{Re}(s)+r>0$, it follows that $D_1(s,r)$ extends analytically to the region $\mathrm{Re}(s)+r>0$ and $\mathrm{Re}(s)>1$ except for a simple pole at $s=1$. Assuming the Riemann Hypothesis, the function may be analytically continued to the half-plane $\mathrm{Re}(s)>-r-\frac{1}{2}$.

\medskip

From Lemma~\ref{dsklemma} and the foregoing discussion, we conclude the following result.

\begin{proposition}\label{arsconvergence}
For a fixed integer $r\geq 1$ and $s=\sigma+it$ with $\sigma,t\in\mathbb{R}$, the Dirichlet series $D_1(s,r)$ has analytic continuation to the half plane for $\mathrm{Re}(s)>-r$ and $\mathrm{Re}(s)>1$, with an exception of a pole at the point $s=1$. Furthermore, this function has analytic continuation to the half-plane $\mathrm{Re}(s)>-r-\frac{1}{2}$ if and only if the Riemann hypothesis is true.    
\end{proposition}

Observe that
\begin{align}\label{cc'relation}
c'_{\chi}(n)
=\sum_{\substack{b\bmod{m}\\(b,m)=1}}\chi(b)e\left(\frac{bn}{m}\right)
=\sum_{\substack{b\bmod{m}}}\chi_0(b)\chi(b)e\left(\frac{bn}{m}\right)
=\sum_{\substack{b\bmod{m}}}\chi_1(b)e\left(\frac{bn}{m}\right)
=c'_{\chi_1}(n),
\end{align}
where $\chi_1$ is a Dirichlet character modulo $m$ defined by $\chi_1=\chi_0\chi$. Note that in Lemma~\ref{shiftedramanujansumformula}, the Dirichlet character $\chi$ is not necessarily a primitive character. Therefore, we further induce it by $\chi^*$ modulo $m^*$, where 
$\chi^*$ is a primitive Dirichlet character such that $m^*|m$.

\footnotetext{Note that $\chi^*_0\chi_0$ is a principal character mod $m$ because $m^*|m$.}

Then we can write $\chi=\chi^*_0\chi^{*}$ for principal character $\chi^*_0$. Then the generalised Gauss sum can be written as
\begin{align}\label{charactersumdecomposition}
c_{{\chi}}(n)
=&\sum_{\substack{b\bmod{m}}}{\chi}_0(b)\chi^{*}(b)e\left(\frac{bn}{m}\right)
=\sum_{\substack{b\bmod{m}\\(b,m)=1}}\chi^{*}(b)e\left(\frac{bn}{m}\right)\nonumber\\
=&\frac{\varphi(m)}{\varphi(m^*)}\sum_{\substack{b\bmod{m^*}\\(b,m^*)=1}}\chi^{*}(b)e\left(\frac{bn}{m^*}\right)
=\frac{\varphi(m)}{\varphi(m^*)}c'_{\chi^*}(n),
\end{align}
where we have used~\cite[Theorem 5.33(a)]{apostol} and~\eqref{generalizedgausssum} in the last step. 

\begin{comment}
Thus, for an arithmetic function $f$, the following identity is true
\begin{align}
\sum_{\chi(\bmod{m})}f(\chi)
=\sum_{m^*|m}\frac{\varphi(m)}{\varphi(m^*)}\sum_{\chi^*(\bmod{m^*})}f(\chi^*).
\end{align}
\end{comment}


\begin{lemma}\label{gausssumlemma}
Let $r>1$ be fixed, and for a Dirichlet character $\chi$ modulo $m$, define
\begin{align}\label{d2sksum}
D_2(s,r)
\coloneqq\sum_{m=1}^{\infty}\frac{1}{\varphi(m)m^{r+1}}\sum_{\substack{\chi(\bmod{m})\\\chi\ne\chi_0}}\tau(\chi)\sum_{n=1}^{\infty}\frac{c'_{\Bar{\chi}}(n)}{n^{s}}.
\end{align}
\noindent
Let $\mathrm{Re}(s)=\sigma$, then the following inequality
\begin{align*}
|D_2(s,r)|
\leq\frac{\zeta(\sigma)\zeta(\sigma+r)\zeta(\sigma+r+1)}{\zeta(2(\sigma+r))}\zeta(r)E_r(\sigma)
\end{align*}
holds for $\sigma>1$, where
\begin{align}\label{erdefinition}
E_r(\sigma)
=C'(r)\prod_{p}(1+J_{p,r}(\sigma))
\end{align}
with
\begin{align*}
J_{p,r}(\sigma)
=-\frac{\frac{1-p^{-r}}{p^{r+1}}\left(\frac{1}{p^{\sigma}}+\frac{1}{2p^{2\sigma+2r+1}}\right)}{1+\frac{1-p^{-r}}{p^{r+1}}}
\end{align*}
and
\begin{align}\label{cprimerdefinition}
C'(r)
\coloneqq\prod_{p}\left(1+\frac{1}{p^{r+1}}\left(1-\frac{1}{p^r}\right)\right).
\end{align}
\end{lemma}

\begin{proof}
Let $\chi$ be a Dirichlet character modulo $m$, induced by a primitive character $\chi^{*}$ modulo $m^{*}$, where $m^{*}|m$. From~\eqref{cc'relation}, it is evident that $c'_{\Bar{\chi}}(n)=c'_{\Bar{\chi}_1}(n)$ with $\bar{\chi}_1=\chi_0\bar{\chi}=\chi_0\chi^*_0\bar{\chi}^*=\chi_0\Bar{\chi}^*$.

\medskip

Beginning with the simplification of the coefficient of the Dirichlet series, and applying Theorem 9.10 and Theorem 9.12 of~\cite{montgomeryvaughanbook}, we obtain

\begin{align*}
%\tau(\chi)c_{\Bar{\chi}}(n)
c'_{\Bar{\chi}}(n)c'_{{\chi}}(1)
=&\chi^{*}\left(\frac{n}{(m,n)}\right)\bar{\chi}^{*}\left(\frac{m/(m,n)}{m^{*}}\right)\mu\left(\frac{m/(m,n)}{m^{*}}\right)\frac{\varphi(m)}{\varphi(m/(m,n))}\mu\left(\frac{m}{m^{*}}\right)\chi^{*}\left(\frac{m}{m^{*}}\right)c'_{\bar{\chi}^{*}}(1)c'_{{\chi^{*}}}(1)\\
=&m^{*}\chi^{*}\left(\frac{n}{(m,n)}\right)\bar{\chi}^{*}\left(\frac{m/(m,n)}{m^{*}}\right)\mu\left(\frac{m/(m,n)}{m^{*}}\right)\frac{\varphi(m)}{\varphi(m/(m,n))}\mu\left(\frac{m}{m^{*}}\right)\chi^{*}\left(\frac{m}{m^{*}}\right).
\end{align*}

\noindent
Let $m=\delta m^*k$ and $n=\delta\ell$ with $(m,n)=\delta$, so that $(m^*k,\ell)=1$. With this change of variables, the above expression becomes
\begin{align}\label{simplifiedcoefficient}
\frac{c'_{\Bar{\chi}}(n)c'_{{\chi}}(1)}{\varphi(m)}
=&m^{*}\chi^{*}\left(\ell\right)\bar{\chi}^{*}\left(k\right)\mu\left(k\right)\frac{1}{\varphi(m^*k)}\mu\left(\delta k\right)\chi^{*}\left(\delta k\right)\nonumber\\
=&m^{*}\chi^{*}\left(\ell\right)\bar{\chi}^{*}\left(k\right)\mu\left(k\right)\frac{1}{\varphi(m^*k)}\mu\left(\delta\right)\mu\left( k\right)\chi^{*}\left(\delta\right)\chi^{*}\left(k\right)\nonumber\\
=&m^{*}\chi^{*}\left(\ell\right)\mu^2\left(k\right)\frac{1}{\varphi(m^*k)}\mu\left(\delta\right)\chi^{*}\left(\delta\right).
\end{align}

%Note that $(\delta,k)=1$ and $(\ell,k)=1$. 
\noindent
Rearranging the terms in~\eqref{d2sksum}, and substituting~\eqref{simplifiedcoefficient} with the change of variable, we obtain
\begin{align*}
D_2(s,r)
=&\sum_{m=1}^{\infty}\frac{1}{m^{r+1}}\sum_{\substack{\chi(\bmod{m})\\\chi\ne\chi_0}}\sum_{n=1}^{\infty}\frac{c'_{\chi}(1)c'_{\Bar{\chi}}(n)}{\varphi(m)n^{s}}\\
%=&\sum_{m=1}^{\infty}\frac{m^*\varphi(m)\mu(m/m^*)}{\varphi(m)m^{r+1}}\sum_{\substack{\chi(\bmod{m})\\\chi\ne\chi_0}}\chi^*(m/m^*)\sum_{n=1}^{\infty}\frac{\chi^*(n/(n,m)m^*)\Bar{\chi}(m/(n,m)m^*)\mu(m/(m,n)m^*)}{\varphi(m/(m,n))n^s}\\
=&\sum_{\delta\geq 1}\mathop{\sum\sum}_{\substack{k,\ell\geq 1\\(\ell,k)=1}}\sum_{m^*\geq 1}\frac{\varphi(m^*k\delta)}{\varphi(m^*)\varphi(m^*k)}\sum_{\chi^*(\bmod{m^*})}\frac{m^*\mu(\delta k)}{(m^*\delta k)^{r+1}}\chi^*(\delta k)\frac{\mu(k)\chi^*(\ell)\Bar{\chi}^*(k)}{(\ell\delta)^s}.
%=&\sum_{k\geq 1}\frac{\mu^2(k)}{k^{r+1}}\sum_{m^*\geq 1}\frac{1}{\varphi(m^{*}){m^*}^{r}}\sum_{\substack{\chi^*(\bmod m^*)}}\sum_{\substack{\delta\geq 1\\(m^*,k)=1\\(\delta,m^*)=1}}\frac{\mu(\delta)\chi^*(\delta)\varphi(\delta)}{\delta^{s+r+1}}\sum_{\substack{\ell\geq 1\\(\ell,m^*k)=1}}\frac{\chi^{*}(\ell)}{\ell^s}.
\end{align*}
\noindent
Furthermore, for $(m^*,k)=1$ and $(\delta,k)=1$, we have
\begin{align*}
\frac{\varphi(m)}{\varphi(m^*k)\varphi(m^*)}
=\frac{\varphi(m^*k\delta)}{\varphi(m^*k)\varphi(m^*)}
=\frac{\varphi(\delta)}{\varphi(m^*)}.
\end{align*}

\noindent
Rearranging the order of summation, we obtain
\begin{align*}
D_2(s,r)
=\sum_{k\geq 1}\frac{\mu^2(k)}{k^{r+1}}\sum_{\substack{m^*\geq 1\\ ({m^*},k)=1}}\frac{1}{\varphi(m^{*}){m^*}^{r}}\sum_{\substack{\chi^*(\bmod m^*)}}\sum_{\substack{\delta\geq 1\\(\delta,k)=1}}\frac{\mu(\delta)\chi^*(\delta)\varphi(\delta)}{\delta^{s+r+1}}\sum_{\substack{\ell\geq 1\\(\ell,k)=1}}\frac{\chi^{*}(\ell)}{\ell^s}.
\end{align*}

\noindent
The sum over $\ell$ equals to
\begin{align}\label{sumoverl}
%\prod_{\substack{p\nmid k\\p\nmid m^*}}\left(1-\frac{\chi^*(p)}{p^s}\right)
L(s,\chi^{*})\prod_{\substack{p\mid k\\p\nmid m^*}}\left(1-\frac{\chi^*(p)}{p^s}\right), 
\end{align}
and the sum over $\delta$ equals to
\begin{align}\label{sumoverdelta}
\prod_{\substack{p\nmid k\\p\nmid m^{*}}}\left(1+\frac{\chi^{*}(p)}{p^{s+r+1}}-\frac{\chi^*(p)}{p^{s+r}}\right).  
\end{align}
\noindent
Observe that, by utilising a decomposition argument analogous to that employed in the preceding lemma, the expression above may be equivalently simplified to
\begin{align*}
\frac{L(s+r+1,\chi^{*})}{L(s+r,\chi^{*})}\prod_{\substack{p|k\\p\nmid m^{*}}}(1+B_{p,r}(s,\chi^*))\left(1+\frac{\chi^{*}(p)}{p^{s+r+1}}-\frac{\chi^*(p)}{p^{s+r}}\right)^{-1},
\end{align*}
where 
\begin{align}\label{bprschidefinition}
B_{p,r}(s,\chi^*)
=-\frac{\left(\frac{{\chi^{*}}(p)^2}{p^{2(s+r)+1}}-\frac{{\chi^{*}}(p)^2}{p^{2(s+r+1)}}+\frac{{\chi^{*}}(p)^3}{p^{3(s+r)+1}}-\frac{{\chi^{*}}(p)^3}{p^{3(s+r+1)}}+\cdots\right)}{\left(1+\frac{\chi^{*}(p)}{p^{s+r+1}}-\frac{\chi^*(p)}{p^{s+r}}-\frac{{\chi^{*}}(p)^2}{p^{2(s+r)+1}}+\frac{{\chi^{*}}(p)^2}{p^{2(s+r+1)}}-\frac{{\chi^{*}}(p)^3}{p^{3(s+r)+1}}+\frac{{\chi^{*}}(p)^3}{p^{3(s+r+1)}}-\cdots\right)}.
\end{align}

\noindent
Note that the stipulation $p\nmid m^*$ in the Euler products mentioned above is superfluous, as $\chi^*(p)=0$ for every prime $p|m^*$. Set
\begin{align}\label{arschidefinition}
A_r(s,\chi^*)
=\prod_{p|k}(1+B_{p,r}(s,\chi^*))\left(1-\frac{\chi^*(p)}{p^s}\right)\left(1+\frac{\chi^{*}(p)}{p^{s+r+1}}-\frac{\chi^*(p)}{p^{s+r}}\right)^{-1}.
\end{align}

\noindent
Thus, the average over the characters is given by
\begin{align}\label{krschidefinition}
K_r(s,\chi^*)
\coloneqq\sum_{\chi^*(\bmod{m^*})}\frac{L(s,\chi^*)L(s+r+1,\chi^*)}{L(s+r,\chi^{*})}A_r(s,\chi^*).
\end{align}

\noindent
We begin by estimating the term $A_r(s,\chi^*)$ and establishing its region of convergence. Letting $s=\sigma+it$ with $\sigma,t\in\mathbb{R}$, and recalling~\eqref{bprschidefinition}, we have
\begin{align*}
|B_{p,r}(s,\chi^*)|
\leq\frac{1}{p^{2\sigma+2r+1}}\frac{1-\frac{1}{2^{\sigma+r+1}}}{1-\frac{1}{2^{\sigma+r}}}
<\frac{1}{2p^{2\sigma+2r+1}},
\end{align*}
and consequently,
\begin{align*}
\sum_{p|k}|B_{p,r}(s,\chi^*)|
<\frac{1}{2}\sum_{p|k}\frac{1}{p^{2\sigma+2r+1}}
<\infty.    
\end{align*}

\noindent
Next, consider the factors
\begin{align*}
\left|1-\frac{\chi^*(p)}{p^{s}}\right|
\leq \left|1+\frac{1}{p^{s}}\right|
=1+\frac{1}{p^{\sigma}}\quad \text{and}\quad &\left|1+\frac{\chi^{*}(p)}{p^{s+r+1}}-\frac{\chi^{*}(p)}{p^{s+r}}\right|
\geq 1-\frac{1}{2^{\sigma+r+1}}.%\quad \text{so that}\\
%&\left|1+\frac{\chi^{*}(p)}{p^{s+r+1}}-\frac{\chi^{*}(p)}{p^{s+r}}\right|^{-1}
%\leq \frac{1}{1-\frac{1}{2^{\sigma+r+1}}}.
\end{align*}

\noindent
It follows that
\begin{align}\label{arskchibound}
|A_r(s,\chi^*)|
\leq&\prod_{p|k}|1+B_{p,r}(s,\chi^*)|\left|1+\frac{1}{p^{s}}\right|\left|1+\frac{1}{p^{s+r+1}}-\frac{1}{p^{s+r}}\right|^{-1}\nonumber\\
\leq&\prod_{p|k}\left(1+\frac{1}{2p^{2\sigma+2r+1}}\right)\left(1+\frac{1}{p^{\sigma}}\right)\left(1-\frac{1}{2^{(\sigma+r+1)}}\right)^{-1}\nonumber\\
\leq&\left(1-\frac{1}{2^{(\sigma+r+1)}}\right)^{-1}\prod_{p|k}\left(1+\frac{1}{p^{\sigma}}+\frac{1}{2p^{2\sigma+2r+1}}+\frac{1}{2p^{3\sigma+2r+1}}\right).
\end{align}

\noindent
Therefore, $|A_r(s,\chi^*)|$ is uniformly bounded in the half-plane $\mathrm{Re}(s)>1$. Moreover, since
\begin{align*}
\frac{\zeta(2\sigma)}{\zeta(\sigma)}
\leq|L(s,\chi^*)|
\leq\zeta(\sigma)    
\end{align*}
for $\mathrm{Re}(s)=\sigma>1$ (see for instance~\cite[Lemma 2.1]{gregmartin}), we may bound the sum in~\eqref{krschidefinition} by
\begin{align*}
|K_r(s,\chi^*)|
\leq&\left|\sum_{\chi^*(\bmod{m^*})}\frac{L(s,\chi^*)L(s+r+1,\chi^*)}{L(s+r,\chi^*)}A_r(s,\chi^*)\right|\\
\leq&\varphi(m^*)\frac{\zeta(\sigma)\zeta(\sigma+r)\zeta(\sigma+r+1)}{\zeta(2(\sigma+r))}\prod_{p|k}\left(1+\frac{1}{p^{\sigma}}+\frac{1}{2p^{2\sigma+2r+1}}\right)\left(1-\frac{1}{2^{\sigma+r+1}}\right)^{-1}.
\end{align*}

\noindent
Ignoring the higher-order terms and inserting the expression above, we obtain
\begin{align*}
|D_2(s,r)|
\leq\frac{\zeta(\sigma)\zeta(\sigma+r)\zeta(\sigma+r+1)}{\zeta(2(\sigma+r))}\sum_{k\geq 1}\frac{\mu^2(k)}{k^{r+1}}\prod_{p|k}\left(1+\frac{1}{p^{\sigma}}+\frac{1}{2p^{2\sigma+2r+1}}\right)\sum_{\substack{m^*\geq 1\\(m^*,k)=1}}\frac{1}{{m^*}^r}.
\end{align*}

\noindent
The sum over $m^*$ equals to
\begin{align*}
\zeta(r)\prod_{p|k}\left(1-\frac{1}{p^{r}}\right).  
\end{align*}

\noindent
Hence, we deduce that
\begin{align*}
|D_2(s,r)|
\leq\frac{\zeta(\sigma)\zeta(\sigma+r)\zeta(\sigma+r+1)\zeta(r)}{\zeta(2(\sigma+r))}\prod_{p}\left(1+\frac{1-p^{-r}}{p^{r+1}}\left(1+\frac{1}{p^{\sigma}}+\frac{1}{2p^{2\sigma+2r+1}}\right)\right),
%=&\frac{\zeta(\sigma)\zeta(\sigma+r+1)}{\zeta(\sigma+r)}E_r(\sigma),
\end{align*}
and set
\begin{align*}
E_r(\sigma)
=\prod_{p}\left(1+\frac{1-p^{-r}}{p^{r+1}}\left(1+\frac{1}{p^{\sigma}}+\frac{1}{2p^{2\sigma+2r+1}}\right)\right).
\end{align*}

\noindent
Furthermore, the Euler product can be expressed in the form
\begin{align*}
E_r(\sigma)
=C'(r)\prod_{p}\frac{\left(1+\frac{1-p^{-r}}{p^{r+1}}(1+E_{p,r}(\sigma))\right)}{1+\frac{1-p^{-r}}{p^{r+1}}},
\end{align*}
where $C'(r)$ is defined in~\eqref{cprimerdefinition} and
\begin{align*}
E_{p,r}(\sigma)
=-\left(\frac{1}{p^{\sigma}}+\frac{1}{2p^{2\sigma+2r+1}}\right).
\end{align*}

\noindent
From our analysis of $B_{p,r}(s,\chi^*)$, it follows that $E_{p,r}(\sigma)$ converges absolutely for $\sigma>1$, thereby completing the proof of the lemma.
\end{proof}


%Note that the growth rates of $D_1(s,r)$ and $D_2(s,r)$ are nearly the same as the upper bound obtained in Lemma~\ref{gausssumlemma}. However, Figure~\ref{figure1} plots the exact values of both sums and illustrates that $D_2(s,r)$ decays more slowly than $D_1(s,r)$ for larger values of $r$ as $\mathrm{Re}(s)$ tends to infinity. 

\subsection{Special values of $D_2(s,r)$}\label{section31}

It is worth noting that the work of Ramar\'{e}~\cite{ramare} can be used to handle the Euler products over arithmetic progressions in~\eqref{sumoverl} and~\eqref{sumoverdelta} and for deriving an exact formula. However, the upper bound established for this Dirichlet series in the preceding lemma suffices for the purposes of this paper. Here, we determine the asymptotic behaviour of $D_2(s,r)$ in the case of a particular Dirichlet character $\chi_4$.



\medskip

Define
\begin{align}\label{chi4definition}
\chi_4(n)
\coloneqq\begin{cases}
1 & \text{if $n\equiv 1(\bmod{4})$}\\
-1 & \text{if $n\equiv 3(\bmod{4})$}\\
0 &  \text{if $2|n$}.
\end{cases}
\end{align}

\noindent
Note that the Dirichlet $L$-function $L(s,\chi_4)=\beta(s)$, the Dirichlet $\beta$-function, is given by
\begin{align*}
L(s,\chi_4)
=\sum_{n=1}^{\infty}\frac{\chi_4(n)}{n^s}
=\sum_{\ell=0}^{\infty}
\frac{(-1)^{\ell}}{(2\ell+1)^s}.
\end{align*}
\noindent
Since $m^*=4$ and $(m^*,k)=1$, it is evident from the proof of Lemma~\ref{gausssumlemma} that
\begin{align*}
D_2(s,r;\chi_4)
=\frac{L(s,\chi_4)L(s+r+1,\chi_4)}{L(s+r,\chi_4)}\sum_{k=1}^{\infty}\frac{\mu^2(k)}{k^{r+1}}A_r(s,\chi_4)
\end{align*}
\noindent
where 
\begin{align*}
A_r(s,\chi_4)
=&\prod_{\substack{p|k}}(1-B_{p,r}(s,\chi_4))\left(1+\frac{\chi_4(p)}{p^{s+r+1}}-\frac{\chi_4(p)}{p^{s+r}}\right)^{-1}\\
=&\prod_{\substack{p|k\\p\equiv 1(\bmod{4})}}(1-B_{p,r}(s))\left(1-\frac{1}{p^s}\right)\left(1+\frac{1}{p^{s+r+1}}-\frac{1}{p^{s+r}}\right)^{-1}\\
&\quad\quad\quad\quad\times\prod_{\substack{p|k\\p\equiv 3(\bmod{4})}}(1-B_{p,r}(s))\left(1+\frac{1}{p^s}\right)\left(1-\frac{1}{p^{s+r+1}}+\frac{1}{p^{s+r}}\right)^{-1},
\end{align*}
\noindent
where $B_{p,r}(s)$ is given in~\eqref{bprschidefinition}. Thus, $D_2(s,r;\chi_4)$ equals to
\begin{align*}
&\frac{L(s,\chi_4)L(s+r+1,\chi_4)}{L(s+r,\chi_4)}\prod_{p\equiv 1(\bmod{4})}\left(1+\frac{(1+p^{-s})}{p^{r+1}}\left(1+\frac{1}{p^{s+r+1}}-\frac{1}{p^{s+r}}\right)^{-1}(1-B_{p,r}(s))\right)\\
&\quad\quad\quad\quad\quad\quad\quad\quad\quad\times\prod_{p\equiv 3(\bmod{4})}\left(1+\frac{(1-p^{-s})}{p^{r+1}}\left(1-\frac{1}{p^{s+r+1}}+\frac{1}{p^{s+r}}\right)^{-1}(1-B_{p,r}(s))\right)\\
=&\frac{\beta(s)\beta(s+r+1)}{\beta(s+r)}\prod_{p\equiv 1(\bmod{4})}\left(1+\frac{(1+p^{-s})}{p^{r+1}}\left(1+\frac{1}{p^{s+r+1}}-\frac{1}{p^{s+r}}\right)^{-1}(1-B_{p,r}(s))\right)\\
&\quad\quad\quad\quad\quad\quad\quad\quad\quad\times\prod_{p\equiv 3(\bmod{4})}\left(1+\frac{(1-p^{-s})}{p^{r+1}}\left(1-\frac{1}{p^{s+r+1}}+\frac{1}{p^{s+r}}\right)^{-1}(1-B_{p,r}(s))\right),
\end{align*}
where $\beta(s)$ denotes the Dirichlet $\beta$ function and we can obtain a decomposed product as before. Note that for $\mathrm{Re}(s)=1$ the above Euler product over arithmetic progression reduced to specific constants.

%As illustrated in Figure~\ref{figure1}, the growth rate of $D_2(s,r)$ is slower than that of $D_1(s,r)$. 

%From the example of $D_2(s,r;\chi_4)$, it is evident that this discrepancy arises due to cancellations occurring in both $D_2(s,r;\chi_4)$ as well as in $D_2(s,r)$, which are induced by the Dirichlet character. Since Lemma~\ref{gausssumlemma} considers only the average order of $D_2(s,r)$, it yields a growth rate comparable to that of $D_1(s,r)$. While an improvement in this estimate is certainly possible, addressing the underlying challenges remains a difficult problem.


\begin{lemma}\label{shfiteddirichletseries}
For a fixed integer $r>1$, the Dirichlet series
\begin{align*}
A_{\sigma_r}(s)
=\sum_{n=1}^{\infty}\frac{\sigma_r(n+1)}{n^s}
=&\zeta(r+1)\sum_{i=0}^{r}\binom{r}{i}D_1(s-i,r)+\zeta(r+1)\sum_{i=0}^{r}\binom{r}{i}D_2(s-i,r),
%=&\zeta(r+1)\sum_{i=0}^{r}\frac{\zeta(s-i)}{\zeta(s+r+1-i)}K_r(s-i)+\zeta(r+1)\sum_{i=0}^{r}
\end{align*}
converges for $\mathrm{Re}(s)>1$, where $D_1(s,r)$ and $D_2(s,r)$ are defined in~\eqref{dskdefinition} and~\eqref{d2sksum}, respectively.
\end{lemma}

\begin{proof}
Employing the identity~\eqref{kthdivisoridentity}, we express
\begin{align*}
A_{\sigma_r}(s)
=\sum_{n=1}^{\infty}\frac{\sigma_r(n+1)}{n^s}
=\zeta(r+1)\sum_{n=1}^{\infty}\frac{(n+1)^{r}}{n^s}\sum_{m=1}^{\infty}\frac{c_m(n+1)}{m^{r+1}}.
\end{align*}
\noindent
Given that $r\geq 1$ is fixed, and under the conditions $\mathrm{Re}(s)>1$, and $\mathrm{Re}(s)+r>1$, one may interchange the order of summation, thereby obtaining
\begin{align}\label{sum1}
\zeta(r+1)\sum_{m=1}^{\infty}\frac{1}{m^{r+1}}\sum_{n=1}^{\infty}\frac{c_m(n+1)}{n^s}(n+1)^r
=\zeta(r+1)\sum_{i=0}^{r}\binom{r}{i}\sum_{m=1}^{\infty}\frac{1}{m^{r+1}}\sum_{n=1}^{\infty}\frac{c_m(n+1)}{n^{s-i}}.
\end{align}

Now applying Lemma~\ref{shiftedramanujansumformula} yields,
\begin{align*}
\sum_{m=1}^{\infty}\frac{1}{m^{r+1}}\sum_{n=1}^{\infty}\frac{c_m(n+1)}{n^{s-i}}
=&\sum_{m=1}^{\infty}\frac{1}{\varphi(m)m^{r+1}}\sum_{n=1}^{\infty}\frac{1}{n^{s-i}}\left(\mu(m)c_m(n)+\sum_{\substack{\chi(\bmod{m})\\ \chi\ne \chi_0}}c_{\chi}(1)c'_{\bar{\chi}}(n)\right)\\
=&\sum_{m=1}^{\infty}\frac{\mu(m)}{\varphi(m)m^{r+1}}\sum_{n=1}^{\infty}\frac{c_m(n)}{n^{s-i}}+\sum_{m=1}^{\infty}\frac{1}{\varphi(m)m^{r+1}}\sum_{\substack{\chi(\bmod{m})\\ \chi\ne \chi_0}}\tau(\chi)\sum_{n=1}^{\infty}\frac{c'_{\bar{\chi}}(n)}{n^{s-i}}\\
=&D_1(s-i,r)+D_2(s-i,r).
\end{align*}

\noindent
By applying Lemma~\ref{dsklemma} and Lemma~\ref{gausssumlemma}, in conjunction with~\eqref{sum1}, we thereby complete the proof.
\end{proof}


We conclude this section by recalling that (see, for instance,~\cite[\S1.3]{titchmarsh}), for $\mathrm{Re}(s)>1$ or $\mathrm{Re}(s)>r+1$,
\begin{align}\label{dsigmars}
D_{\sigma_r}(s)
=\sum_{n=1}^{\infty}\frac{\sigma_r(s)}{n^s}
=\zeta(s)\zeta(s-r),
\end{align}
for fixed $r$.


\section{The Integral Asymptotics}

As outlined in Section~\ref{initialsetup}, we have two main steps to establish the central limit theorem. First, to employ the saddle-point method, we analyse the asymptotic behaviour of the integrands

\begin{align}\label{maintermintegral}
\Omega_{1,f}(n)
\coloneqq\frac{e^{n\tau}}{2\pi}\int_{|\theta|\leq \theta_n}\exp(in\theta+F(\tau+i\theta,u))d\theta,
\end{align}
and
\begin{align}\label{errortermintegral}
\Omega_{2,f}(n)
\coloneqq\frac{e^{n\tau}}{2\pi}\int_{\theta_n<|\theta|\leq \pi}\exp(in\theta+F(\tau+i\theta,u))d\theta,
\end{align}
with a suitable choice of $\tau=\tau(n,u)$ such that the first-order derivative in~\eqref{taylorapproximation} vanishes.


\subsection{Estimate of $\Omega_{1,f}$}\label{EstimateofOmega1f}
%We study the asymptotic behaviour of $\Omega_{1,f}(n)$ to solve for the saddle points by utilizing the Taylor expansion provided in equation~\eqref{taylorapproximation}. 
From the definition of $F(\gamma,u)$,  it is evident that deriving an asymptotic expression for the integrand $\Omega_{1,f}$ necessitates a comprehensive understanding of the analytic properties of the Dirichlet series with shifted coefficient $A_f(s)$ and $D_f(s)$ corresponding to $f(n)=\sigma_r(n)$.

\begin{comment}
Note that since $\chi$ is a primitive Dirichlet character $\chi(-1)=\pm 1$, for even and odd characters. We then simplify the inner sum of~\eqref{generalizedramanujansum} as
\begin{align*}
c_m(n+1)
=&\frac{m}{\varphi(m)}\sum_{\substack{\chi(\bmod{m})}}{\chi}(-1)\chi(n),\quad\text{if $(m,n)=1$}\\
=&\frac{m}{2\varphi(m)}\left(\sum_{\substack{\chi(\bmod{m})\\{\chi\mathrm{ even}}}}\chi(n)-\sum_{\substack{\chi(\bmod{m})\\\chi\mathrm{ odd}}}\chi(n)\right)\quad\text{if $(m,n)=1$}.
\end{align*}

Note that for primitive Dirichlet chacarter $\chi$,
\begin{align*}
\sum_{\substack{\chi(\bmod{m})\\{\chi\;\mathrm{even}}}}\chi(n)  =\frac{1}{2}\sum_{\substack{\chi(\bmod{m})}}(\chi(r)+\chi(-r)),\quad\text{and}
\sum_{\substack{\chi(\bmod{m})\\{\chi\;\mathrm{odd}}}}\chi(n)  =\frac{1}{2}\sum_{\substack{\chi(\bmod{m})}}(\chi(r)-\chi(-r)).
\end{align*}

Since $(m,n)=1$, from~\cite[Lemma 4.1]{buietal}, we use the identity
\begin{align*}
\sum_{\substack{\chi(\bmod{m})\\{\chi\mathrm{ even}}}}\chi(n)  =\frac{1}{2}\sum_{\substack{dr=m\\d|n\pm 1}}\mu(r)\varphi(d),  
\end{align*}
to conclude the lemma.

    

\begin{lemma}\label{maintermforomega1qu}
Let $c_m(n)$ be the Ramanujan sum as given in~\eqref{ramanujansum}. Set
\begin{align*}
\mathcal{M}_{1,j}(s,u)
\coloneqq\esum_{\chi(\bmod{m})}L(s-j,\chi)\mathrm{Li}_{s-j+1}(-u)\Gamma(s).
\end{align*}

Then 
\begin{align*}
\mathcal{M}_{1,j}(s,u)
=
\end{align*}


\end{lemma}

\begin{proof}
It is worth noting that the proceeding argument is analogous to the analysis of the first moment of the Dirichlet $L$-function. The functional equation of $L(s,\chi)$ can be expressed as

\begin{align*}
\Lambda(s,\chi)=\varepsilon(\chi)\Lambda(1-s,\Bar{\chi})
=\left(\frac{\pi}{m}\right)^{-\frac{s+a}{2}}\Gamma\left(\frac{s+a}{2}\right)L(s,\chi),
\end{align*}
where $|\varepsilon(\chi)|=1$ is the root number of the $L$-function and
\begin{align*}
a=a(\chi)=
\begin{cases}
0 & \text{if}\; \chi(-1)=1,\\   
1 & \text{if}\; \chi(-1)=-1.
\end{cases}
\end{align*}

Note that $\Lambda(s,\chi)$ has analytic continuation to the entire complex plane. Let
\begin{align*}
I_{1}(m,\gamma,u)
=\frac{1}{2\pi i}\int_{(c)}\Lambda(s,\chi)\frac{G(s)}{s}ds\;
=\frac{1}{2\pi i}\int_{(c)}\left(\frac{\pi}{m}\right)^{-s/2}\Gamma\left(\frac{s}{2}\right)L(s,\chi)\frac{G(s)}{s}ds,  \end{align*}
where $G(s)$ is even and $G(0)=1$. Now, we change the variable from $s\to -s$ and use the functional equation again to get
\begin{align*}
I_{2}(m,\gamma,u)
=-\frac{1}{2\pi i}\int_{(c)}\Lambda(1-s,\chi)\frac{G(s)}{s}ds\;
=-\frac{1}{2\pi i}\int_{(c)}\left(\frac{\pi}{m}\right)^{-1+s/2}\Gamma\left(1-\frac{s}{2}\right)L(1-s,\chi)\frac{G(s)}{s}ds, 
\end{align*}

By the approximate functional equation of the $L$-functions (see Theorem 5.3 of~\cite{iwanieckowlaski}) we have

\begin{align*}
L(s,\chi)
=\sum_{n=1}^{\infty}\frac{\chi(n)}{n^s}V_s\left(\frac{n}{\sqrt{m}/\pi}\right)+\varepsilon(\chi)\sum_{n=1}^{\infty}\frac{\Bar{\chi(n)}}{n^{1-s}}V_{1-s}\left(\frac{n}{\sqrt{m}\pi}\right),
\end{align*}

where
\begin{align*}
V_t(x)=\frac{1}{2\pi i}\int_{(1)}\frac{G(t)}{t}\left(\frac{\pi}{m}\right)^{-\frac{t+a}{2}}\Gamma\left(\frac{t+a}{2}\right)x^{-t} dt, 
\end{align*}

and $\varepsilon(\chi)=\tau(\chi)/\sqrt{m}$.
\end{proof}



%\subsubsection{Estimate of $\Omega_{1,\sigma_k}$} 

%In this section, we establish a corresponding result to Lemma~\ref{maintermlemmaforcqn} for the divisor function $\sigma_k(n)$, where $k \geq 1$ is fixed. The key difficulty in the preceding argument arises from the fact that the moduli $m$ in the Ramanujan sum is not fixed (see, for instance,~\eqref{kthdivisoridentity}).
\end{comment}




%We are now ready to compute the main term.

\begin{lemma}\label{maintermforsigmar}
Let $f$ denote the $r$-th divisor function $\sigma_r(n)$ for fixed $r\geq 2$, and let $B^2=F_{\gamma\gamma}(\tau,u)$. We then have
\begin{align*}
\Omega_{1,f}
=\frac{e^{n\tau+F(\tau, u)}}{2\pi}\left(\int_{-\theta_n}^{\theta_n}\exp\left(-\frac{B^2}{2}\theta^2\right)d\theta\right)\left(1+O\left(\tau^{2r/7-1}\right)\right),
\end{align*}
where
\begin{align*}
\int_{-\theta_n}^{\theta_n}\exp\left(-\frac{B^2}{2}\theta^2\right)d\theta 
=\frac{\sqrt{2\pi}}{B}+O\left(\tau^{-1-3r/7}B^{-2}\exp\left(-\frac{\tau^{-r/7-1}}{2}\right)\right)
\end{align*}
uniformly in $u$. 
\end{lemma}


\begin{proof}
Note that
\begin{align}\label{saddlepointequationforsigmak}
F(\tau+i\theta,u)
=F(\tau,u)+F_{\gamma}(\tau, u)i\theta-\frac{F_{\gamma\gamma}(\tau,u)}{2}\theta^2+O\left(\theta^3\sup_{0\leq \theta_0\leq \theta}|F_{\gamma\gamma\gamma}(\tau+i\theta_0,u)|\right)
\end{align}

\noindent
Utilising the definition of $\Delta_f(n)$ (as given in~\eqref{gfunctiondefintion}), we subsequently deduce that

\begin{align}\label{curlyfgammau}
F(\gamma,u)
=&\sum_{n=1}^{\infty}\Delta_f(n)\log(1+ue^{-n\gamma})
=-\sum_{n=1}^{\infty}\Delta_f(n)\sum_{\ell=1}^{\infty}\frac{(-u)^{\ell}}{\ell}e^{-n\ell\gamma}\nonumber\\
=&-\sum_{n=1}^{\infty}(\sigma_r(n+1)-\sigma_r(n))\sum_{\ell=1}^{\infty}\frac{(-u)^{\ell}}{\ell}e^{-n\ell\gamma}\nonumber\\
=&-\sum_{n=1}^{\infty}\sigma_r(n+1)\sum_{\ell=1}^{\infty}\frac{(-u)^{\ell}}{\ell}e^{-n\ell\gamma}
+\sum_{n=1}^{\infty}\sigma_r(n)\sum_{\ell=1}^{\infty}\frac{(-u)^{\ell}}{\ell}e^{-n\ell\gamma}\nonumber\\
%&\quad\quad\quad\quad\quad\quad\quad\quad\quad\quad\quad\quad\quad\quad-\sum_{n=1}^{\infty}\sigma_k(n+1)\sum_{\ell=1}^{\infty}\frac{(-u)^{\ell}}{\ell}e^{-n\ell\gamma}\nonumber\\
=&-\mathcal{F}_{1}(\gamma,u)+\mathcal{F}_{2}(\gamma,u).
%-\mathcal{F}_{3}(\gamma,u).
\end{align}

\medskip

We begin with the case $\mathcal{F}_{2}(\gamma,u)$ owing to its relative simplicity. For $j\geq 0$, we have
\begin{align}\label{curlyf2gammau}
\frac{\partial^j}{\partial\gamma^j}\mathcal{F}_{2}(\gamma,u)
=(-1)^j\sum_{n=1}^{\infty}\sigma_r(n)n^{j}\sum_{\ell=1}^{\infty}(-u)^{\ell}\ell^{j-1}e^{-n\gamma\ell}.
\end{align}
\noindent
Observe that all the infinite sums above take the form
\begin{align*}
\mathfrak{h}_{2,j}(\gamma,u)
=\sum_{n=1}^{\infty}\sigma_r(n)n^{j}\sum_{\ell=1}^{\infty}(-u)^{\ell}\ell^{j-1}e^{-n\gamma\ell}.
\end{align*}
\noindent
Let $\mathcal{H}_{2,j}(s,u)$ denote the Mellin transform of $\mathfrak{h}_{2,j}(\gamma,u)$ with respect to $\gamma$. Therefore, by~\eqref{dsigmars}
\begin{align*}
\mathcal{H}_{2,j}(s,u)
=&\sum_{n=1}^{\infty}\sigma_r(n)n^{j}\sum_{\ell=1}^{\infty}(-u)^{\ell}\ell^{j-1}\int_{0}^{\infty}e^{-n\ell\gamma}\gamma^{s-1}d\gamma\\
=&\sum_{n=1}^{\infty}\sigma_r(n)n^{j}\sum_{\ell=1}^{\infty}(-u)^{\ell}\ell^{j-1}\Gamma(s)(n\ell)^{-s}\\
=&\sum_{n=1}^{\infty}\frac{\sigma_r(n)}{n^{s-j}}\sum_{\ell=1}^{\infty}\frac{(-u)^{\ell}}{\ell^{s-j+1}}\Gamma(s)
=\zeta(s-j)\zeta(s-j-r)\mathrm{Li}_{s-j+1}(-u)\Gamma(s).
\end{align*}
\noindent
Note that $\mathcal{H}_{2,j}(s,u)$ converges for $\mathrm{Re}(s)>r+j+2$ and $\mathrm{Re}(s)>j+2$, with its double poles located at the intervals $\frac{3}{2}+j\leq \mathrm{Re}(s)\leq j+2$ and $\frac{3}{2}+r+j\leq\mathrm{Re}(s)\leq r+j+2$ for fixed $r\geq 2$. These poles occur at $s=j+1$ and $s=j+r+1$. Consequently, it is necessary to invoke Theorem~\ref{cmtheorem} twice. Thus,
\begin{align*}
\mathfrak{h}_{2,j}(\gamma,u)
=\zeta(r+1)\mathrm{Li}_{r+2}(-u)\Gamma(j+r+1)\gamma^{-(j+r+1)}&+\zeta(1-r)\mathrm{Li}_{2}(-u)\Gamma(j+1)\gamma^{-(j+1)}\\
&\quad\quad\quad\quad+O(\gamma^{-(j+3/2)}).
\end{align*}

\noindent
Thus, inserting the above expression in~\eqref{curlyf2gammau}, we arrive at
\begin{align}\label{f2derivativefinal}
\frac{\partial^j}{\partial\gamma^j}\mathcal{F}_{2}(\gamma,u)
=(-1)^j\mathfrak{h}_{2,j}(\gamma,u)
=&(-1)^j\zeta(r+1)\mathrm{Li}_{r+2}(-u)\Gamma(j+r+1)\gamma^{-(j+r+1)}\nonumber\\
&+(-1)^j\zeta(1-r)\mathrm{Li}_{2}(-u)\Gamma(j+1)\gamma^{-(j+1)}+O(\gamma^{-(j+r+3/2)})\nonumber\\
=&(-1)^{j}\zeta(r+1)\mathrm{Li}_{r+2}(-u)\Gamma(j+r+1)\gamma^{-(j+r+1)}+O(\gamma^{-(j+r+3/2)}).
\end{align}

\noindent
We now proceed with the computation of $\mathcal{F}_{1}(\gamma,u)$. For $j\geq 0$, we write
\begin{align}\label{f1derivative}
\frac{\partial^j}{\partial\gamma^j}\mathcal{F}_{1}(\gamma,u)
=(-1)^j\sum_{n=1}^{\infty}\sigma_r(n+1)n^{j}\sum_{\ell=1}^{\infty}(-u)^{\ell}\ell^{j-1}e^{-n\ell\gamma}.
\end{align}

\noindent
Similarly, as before, all the infinite sums above take the form
\begin{align*}
\mathfrak{h}_{1,j}(\gamma,u)
=\sum_{n=1}^{\infty}\sigma_r(n+1)n^{j}\sum_{\ell=1}^{\infty}(-u)^{\ell}\ell^{j-1}e^{-n\ell\gamma}.
\end{align*}

\noindent
Let $\mathcal{H}_{1,j}(s,u)$ denote the Mellin transform of $\mathfrak{h}_{1,j}(\gamma,u)$ with respect to $\gamma$. Thus, we have
\begin{align*}
\mathcal{H}_{1,j}(s,u)
=&\sum_{n=1}^{\infty}\sigma_r(n+1)n^j\sum_{\ell=1}^{\infty}(-u)^{\ell}\ell^{j-1}\int_{0}^{\infty}e^{-\ell\gamma}\gamma^{s-1}d\gamma\nonumber\\
=&\sum_{n=1}^{\infty}\sigma_r(n+1)n^j\sum_{\ell=1}^{\infty}(-u)^{\ell}\ell^{j-1}\Gamma(s)(n\ell)^{-s}\nonumber\\
=&\sum_{n=1}^{\infty}\frac{\sigma_r(n+1)}{n^{s-j}}\sum_{\ell=1}^{\infty}\frac{(-u)^{\ell}}{\ell^{s-j+1}}\Gamma(s).
\end{align*}

\noindent
Applying Lemma~\ref{shfiteddirichletseries}, the aforementioned expression becomes
\begin{align}\label{h1jmellintransform}
\mathcal{H}_{1,j}(s,u)
=&\zeta(r+1)\left(\sum_{i=0}^{r}\binom{r}{i}D_1(s-i-j,r)+\sum_{i=0}^{r}\binom{r}{i}D_2(s-i-j,r)\right)\mathrm{Li}_{s-j+1}(-u)\Gamma(s).
%=&\widetilde{\mathfrak{h}}_{1,j}(\gamma,u)+\doublewidetilde{\mathfrak{h}}_{1,j}.
\end{align}

\noindent
Recalling Lemma~\ref{dsklemma}, the first term in the above expression yields
\begin{align*}
\widetilde{\mathcal{H}}_{1,j}(s,u)
=&\zeta(r+1)\sum_{i=0}^{r}\binom{r}{i}D_1(s-i-j,r)\mathrm{Li}_{s-j+1}(-u)\Gamma(s)\\
=&\zeta(r+1)\sum_{i=0}^{r}\binom{r}{i}\frac{\zeta(s-i-j)}{\zeta(s-i-j+r+1)}K_{r}(s-i-j)\mathrm{Li}_{s-j+1}(-u)\Gamma(s),
\end{align*}
where $K_r(s)$ is defined in~\eqref{arsdefinition}. Notably, $K_r(s)$ is absolutely convergent for $\mathrm{Re}(s)-i-j+r>0$. The only singularity arises from the zeta function at $s=1+i+j$. By applying Theorem~\ref{cmtheorem}, we obtain
\begin{align*}
\widetilde{\mathfrak{h}}_{1,j}(\gamma,u)
=&\frac{\zeta(r+1)K_r(1)}{\zeta(r+2)}\sum_{i=0}^{r}\binom{r}{i}\gamma^{-(i+j+1)}\mathrm{Li}_{i+2}(-u)\Gamma(i+j+1)+O(\gamma^{-r-j-3/2}).
\end{align*}

\noindent
By Lemma~\ref{gausssumlemma}, we obtain an upper bound for the second term in~\eqref{h1jmellintransform} as follows
\begin{align*}
|\doublewidetilde{\mathcal{H}}_{1,j}(s,u)|
=&\zeta(r+1)\left|\sum_{i=0}^{r}\binom{r}{i}D_2(s-i-j,r)\right|
\leq\zeta(r+1)\sum_{i=0}^{r}\binom{r}{i}|D_2(s-i-j,r)|\\
\leq&\zeta(r+1)\zeta(r)\sum_{i=0}^{r}\binom{r}{i}\frac{\zeta(\sigma-i-j)\zeta(\sigma-i-j+r)\zeta(\sigma-i-j+r+1)}{\zeta(2(\sigma-i-j+r))}\\
&\quad\quad\quad\quad\quad\quad\quad\quad\quad\quad\quad\quad\quad\quad\quad\quad\quad\quad E_r(\sigma-i-j)\mathrm{Li}_{s-j+1}(-u)\Gamma(s),
\end{align*}
where $\mathrm{Re}(s)=\sigma$ and $E_r(\sigma)$ is defined in~\eqref{erdefinition}. However, triple poles arise from the zeta functions at $\sigma=i+j+1$, $\sigma=i+j+1-r$, and $\sigma=i+j-r$. By invoking Theorem~\ref{cmtheorem}, we express this as
\begin{align*}
|\doublewidetilde{\mathfrak{h}}_{1,j}(\gamma,u)|
\leq&\zeta(r+1)\zeta(r)\sum_{i=0}^{r}\binom{r}{i}\frac{\zeta(r+1)\zeta(r+2)}{\zeta(2(r+1))}E_r(1)\mathrm{Li}_{i+2}(-u)\Gamma(i+j+1)\gamma^{-(i+j+1)}\\
&+\zeta(r+1)\zeta(r)\sum_{i=0}^{r}\binom{r}{i}\zeta(1-r)E_r(1-r)\mathrm{Li}_{i+2-r}(-u)\Gamma(i+j+1-r)\gamma^{-(i+j+1-r)}\\
&+\zeta(r+1)\zeta(r)\sum_{i=0}^{r}\binom{r}{i}{\zeta(-r)E_r(-r)}\mathrm{Li}_{i+1-r}(-u)\Gamma(i+j-r)\gamma^{-i-j+r}+O(\gamma^{-r-j-3/2})\\
=&\frac{\zeta^2(r+1)\zeta(r+2)\zeta(r)}{\zeta(2(r+1))}E_r(1)\sum_{i=0}^{r}\binom{r}{i}\mathrm{Li}_{i+2}(-u)\Gamma(i+j+1)\gamma^{-(i+j+1)}+O(\gamma^{-j-r-3/2}).
\end{align*}

\footnotetext{Note that due to the convergence region of $E_r(\sigma)$, $E_r(-r)$ and $E_r(1-r)$ are constants.}
\noindent
Note that the average orders of $\widetilde{\mathfrak{h}}_{1,j}(\gamma,u)$ and $\doublewidetilde{\mathfrak{h}}_{1,j}(\gamma,u)$ are of the same size, as established in Lemma~\ref{dsklemma} and Lemma~\ref{gausssumlemma}. Substituting these expressions into~\eqref{f1derivative}, we obtain
\begin{align}\label{f1derivativefinal}
&\frac{\partial^j}{\partial\gamma^j}\mathcal{F}_1(\gamma,u)
=(-1)^{j}\mathfrak{h}_{1,j}(\gamma,u)
=(-1)^{j}\;\widetilde{\mathfrak{h}}_{1,j}(\gamma,u)+(-1)^{j}\;\doublewidetilde{\mathfrak{h}}_{1,j}(\gamma,u)\nonumber\\
\leq&\frac{\zeta(r+1)K_r(1)}{\zeta(r+2)}\sum_{i=0}^{r}\binom{r}{i}\gamma^{-(i+j+1)}\mathrm{Li}_{i+2}(-u)\Gamma(i+j+1)\nonumber\\
&+\frac{\zeta^2(r+1)\zeta(r+2)\zeta(r)}{\zeta(2(r+1))}E_r(1)\sum_{i=0}^{r}\binom{r}{i}\gamma^{-(i+j+1)}\mathrm{Li}_{i+2}(-u)\Gamma(i+j+1)+O(\gamma^{-j-r-3/2})\nonumber\\
&\leq\left(\frac{\zeta(r+1)K_r(1)}{\zeta(r+2)}+\frac{\zeta^2(r+1)\zeta(r+2)}{\zeta(2(r+1))}\zeta(r)E_r(1)\right)\sum_{i=0}^{r}\binom{r}{i}\gamma^{-(i+j+1)}\mathrm{Li}_{i+2}(-u)\Gamma(i+j+1)\nonumber\\
&\quad\quad\quad\quad\quad\quad\quad\quad\quad\quad\quad\quad\quad\quad\quad\quad\quad\quad\quad\quad\quad\quad\quad\quad+O(\gamma^{-j-r-3/2}).
\end{align}

\noindent
Thus, by combining equations~\eqref{f1derivativefinal} and~\eqref{f2derivativefinal}, we deduce that~\eqref{curlyfgammau} becomes
\begin{align}\label{fgammauderivativefinal}
&\frac{\partial^j}{\partial\gamma^j}F(\gamma,u)
=-\frac{\partial^j}{\partial\gamma^j}\mathcal{F}_1(\gamma,u)
+\frac{\partial^j}{\partial\gamma^j}\mathcal{F}_2(\gamma,u)
=(-1)^{j+1}\mathfrak{h}_{1,j}(\gamma,u)+(-1)^{j}\mathfrak{h}_{2,j}(\gamma,u)\nonumber\\
\leq &\bigg|\left(\frac{\zeta(r+1)K_r(1)}{\zeta(r+2)}+\frac{\zeta^2(r+1)\zeta(r+2)}{\zeta(2(r+1))}\zeta(r)E_r(1)\right)\sum_{i=0}^{r}\binom{r}{i}\gamma^{-(i+j+1)}\mathrm{Li}_{i+2}(-u)\Gamma(i+j+1)\nonumber\\
&-\zeta(r+1)\mathrm{Li}_{r+2}(-u)\Gamma(r+j+1)\gamma^{-(r+j+1)}\bigg|+O(\gamma^{-j-r-3/2})\nonumber\\
\leq& N(r)\mathrm{Li}_{r+2}(-u)\Gamma(r+j+1)\gamma^{-(r+j+1)}+O(\gamma^{-j-r-3/2}),
\end{align}
where
\begin{align}\label{nrdefinition}
N(r)
\coloneqq\left|\frac{\zeta(r+1)K_r(1)}{\zeta(r+2)}+\frac{\zeta^2(r+1)\zeta(r+2)}{\zeta(2(r+1))}\zeta(r)E_r(1)-\zeta(r+1)\right|.  
\end{align}
\noindent
It is important to note that due to the oscillatory behaviour of the gap function $\Delta_f(n)$ for $f(n)=\sigma_r(n)$, together with the fact that we can only derive an upper bound for $D_2(s,r)$, we obtain an average order for the derivative of $F(\gamma,u)$ rather than an exact formula.

%Observe that the terms in the final two lines of~\eqref{fgammauderivativefinal} will be absorbed into the error term, as their order of magnitude is lower than that of the main term and hence negligible in comparison. 

\medskip

Thus, for the saddle-point $n$, we have
\begin{align}\label{saddlepointsolutionforsigmak}
n=-F_{\gamma}(\tau,u)
=&O\left(N(r)\mathrm{Li}_{r+2}(-u)\Gamma(r+2)\tau^{-(r+2)}\right)
%&\bigg(1+\sum_{i=0}^{r-1}\mathrm{Li}_{i+2}(-u)\Gamma(i+2)\tau^{-i}\bigg)-1\bigg)+O\left(\tau^{-2}\right),
\end{align}
whereas
\begin{align*}
B^2=F_{\gamma\gamma}(\tau,u)
=&O\left(N(r)\mathrm{Li}_{r+2}(-u)\Gamma(r+3)\tau^{-(r+3)}\right).
\end{align*}

\noindent
For the third derivative appearing in the error term of~\eqref{saddlepointequationforsigmak}, we have
\begin{align*}
\frac{\partial^3}{\partial \theta^3}F(\tau+i\theta,u)
=-i\sum_{k=1}^{\infty}\frac{k^3\Delta_f(k)(1-u^{-1}e^{k(\tau+i\theta)})}{(1+u^{-1}e^{k(\tau+i\theta)})^3}.
\end{align*}

\noindent
The following estimation technique is adapted from~\cite{madwag}. Let $k_0=\tau^{-(1+c)}$ for some $c>0$, and for simplicity of notation, set $v=u^{-1}$. We consider two cases: first, when $k$ is large i.e., $k>k_0$; and second, when $k\leq k_0$. Beginning with the large sum, we obtain
\begin{align}\label{thirdorderderivative}
&\left|-i\sum_{k>k_0}\frac{k^3\Delta_f(k)(1-ve^{k(\tau+i\theta)})}{(1+ve^{k(\tau+i\theta)})^3}\right|
\leq\sum_{k>k_0}\frac{k^3|\Delta_f(k)|(1+ve^{k\tau})}{|1+ve^{k(\tau+i\theta)}|^3}\nonumber\\
\leq&\sum_{k>k_0}\frac{k^3|\Delta_f(k)|(1+ve^{k\tau})}{(ve^{k\tau}-1)^3}
=O\left(\sum_{k>k_0}\frac{k^3|\Delta_f(k)|}{e^{k\tau}}\right)
=O\left(\tau^{-r-4}\right),
\end{align}
where we have applied the upper bound (see, for instance,~\cite{ramanujan})
\begin{align*}
\sigma_r(n)
<n^r(1+2^{-r}+3^{-r}+\cdots)
<n^r\zeta(r).    
\end{align*}

\noindent
For the remaining case, we observe that
\begin{align*}
\left|1+v^{k(\tau+i\theta)}\right|
\geq (1+ve^{k\tau})\cos{\left(\frac{k\theta}{2}\right)}.
\end{align*}

\noindent
Thus, we obtain
\begin{align*}
&\left|-i\sum_{k\leq k_0}\frac{k^3\Delta_f(k)(1-u^{-1}e^{k(\tau+i\theta)})}{(1+u^{-1}e^{k(\tau+i\theta)})^3}\right| 
\leq\sum_{k\leq k_0}\frac{k^3|\Delta_f(k)|(1+ve^{k\tau})}{|1+ve^{k(\tau+i\theta)}|^3}\\
\leq&\sum_{k\leq k_0}\frac{k^3|\Delta_f(k)|(1+ve^{k\tau})}{(1+ve^{k(\tau+i\theta)})^3}\left(1+O((k\theta)^2)\right)
\leq \sum_{k\leq k_0}\frac{k^3|\Delta_f(k)|}{ve^{k\tau}}\left(1+O((k\theta)^2)\right)\\
\leq&\sum_{k\geq 1}\frac{k^3|\Delta_f(k)|}{ve^{k\tau}}+O\left(\sum_{k\geq 1}\frac{k^5|\Delta_f(k)|\theta^2}{e^{k\tau}}\right).
\end{align*}

\noindent
Applying the Mellin transform and converse mapping, we have
\begin{align*}
\left|-i\sum_{k=1}^{\infty}\frac{k^3\Delta_f(k)(1-u^{-1}e^{k(\tau+i\theta)})}{(1+u^{-1}e^{k(\tau+i\theta)})^3}\right|
=O\left(\tau^{-r-4}+\tau^{-r-6}\theta^2\right).
\end{align*}

\noindent
Combining this with~\eqref{thirdorderderivative}, we get
\begin{align*}
\frac{\partial^3}{\partial\theta^3}F(\tau+i\theta,u)
=O(\tau^{-r-4}),
\end{align*}
for $|\theta|\leq\tau^{1+3r/7}$. All these together leads to the expansion
\begin{align*}
F(\tau+i\theta,u)
=F(\tau,u)+in\theta-\frac{B^2}{2}\theta^2+O(\tau^{-r-4}\theta^3).
\end{align*}

\noindent
Thus for the integral $\Omega_{1,f}$, we obtain
\begin{align*}
&\frac{e^{n\tau}}{2\pi}\int_{-\theta_n}^{\theta_n}\exp(in\theta+F(\tau+i\theta,u))d\theta\\
=&\frac{e^{n\tau+F(\tau,u)}}{2\pi}\int_{-\theta_n}^{\theta_n}\exp\left(-\frac{B^2}{2}\theta^2+O(\tau^{2r/7-1})\right)d\theta\\
=&\frac{e^{n\tau+F(\tau,u)}}{2\pi}\left(\int_{-\theta_n}^{\theta_n}\exp\left(-\frac{B^2}{2}\theta^2\right)d\theta\right)\left(1+O(\tau^{2r/7-1})\right).
\end{align*}

\noindent
Finally, we transform the integral into a Gaussian integral and arrive at
\begin{align*}
\int_{-\theta_n}^{\theta_n}\exp\left(-\frac{B^2}{2}\theta^2\right)d\theta
=&\int_{-\infty}^{\infty}\exp\left(-\frac{B^2}{2}\theta^2\right)d\theta
-2\int_{\theta_n}^{\infty}\exp\left(-\frac{B^2}{2}\theta^2\right)d\theta\\
=&\frac{\sqrt{2\pi}}{B}+O\left(\int_{\theta_n}^{\infty}\exp\left(-\frac{B^2\tau^{1+3r/7}}{2}\theta\right)d\theta\right)\\
=&\frac{\sqrt{2\pi}}{B}+O\left(\tau^{-1-3r/7}B^{-2}\exp\left(-\frac{\tau^{-r/7-1}}{2}\right)\right),
\end{align*}
completing the proof.
\end{proof}


\subsection{Estimate of $\Omega_{2,f}$}


%\subsubsection{Estimate of $\Omega_{2,\sigma_k}$}

We proceed to prove the following lemma, which establishes an upper bound for the integrand in~\eqref{errortermintegral} when $f(n)=\sigma_r(n)$ for fixed $r\geq 2$. 

\begin{lemma}\label{omega2sigmak}
For a fixed $r\geq 2$ with $f(n)=\sigma_r(n)$, the error term is bounded as follows
\begin{align*}
|\Omega_{2,\sigma_r}|
=O\left(\exp\left(n\tau+F(\tau,u)-c_{r}\tau^{-r/7}\right)\right),    
\end{align*} 
where $c_{r}$ is a constant uniformly in $u$.
\end{lemma}

We first note the elementary inequality $\sigma_r(n)>n^r$ for $n\geq 2$ and $r\geq 1$,  which simplifies the analysis in this section. Additionally, we incorporate the following lemma by Li and Chen~\cite[Lemma 2.5]{lichen} as a key element of the preceding argument. 

\begin{lemma}[Li-Chen,~\cite{lichen}]\label{lichenlemma}
Let $k\geq 1$ be fixed. Then for $\xi>0$, we have
\begin{align*}
\sum_{n\geq 1}n^{k-1}e^{-n\xi}(1-\cos(ny))\geq \varrho\left(\frac{e^{-\xi}}{(1-e^{-\xi})^{k}}-\frac{e^{-\xi}}{\left|1-e^{-\xi-iy}\right|^{k}}\right),   
\end{align*}
where $\varrho$ is a positive constant which depends only on $k$.
\end{lemma}

We are now ready to prove the following result, which will directly imply Lemma~\ref{omega2sigmak}. In this context, it is noteworthy that, owing to the inequality $\sigma_r(n)>n^r$, for $n\geq 2$ and $r\geq 1$, the following lemma bears resemblance to Lemma 7 in~\cite{lipniketal}. %However, we provide the argument in full detail for the convenience of the readers.

\begin{lemma}\label{qlemmaforomega2sigmak}
For a fixed $r\geq 2$ and any real $y$ satisfying $\theta_0\leq |y|\leq \pi$, we have
\begin{align*}
\frac{|Q_{\sigma_r}(e^{\tau+iy},u)|}{Q_{\sigma_r}(e^{-\tau,u})}
\leq\left(\exp\left(-\frac{2u(r-1)}{(1+u)^2}\varrho\left(\frac{1}{4}\tau^{-{r}/{7}}\right)\right)\right),
\end{align*}
where $\varrho$ is a constant depends on $r$.
\end{lemma}

\begin{proof}
Utilising the inequality $\sigma_r(n)>n^r$ for $n\geq 2$ and $r\geq 2$, we can simplify
\begin{align*}
|\Delta_f(n)|
=|\sigma_r(n+1)-\sigma_r(n)|
\geq c((n+1)^{r}-n^{r})
=c\sum_{i=0}^{r-1}\binom{r}{i}n^{i},
\end{align*}
for some constant $c\geq 0$. However, for $c=0$ in the above inequality, the following argument becomes trivial. So, we only consider the case where $c>0$.

By the mean value theorem, there exists a $\xi\in (n,n+1)$ such that $(n+1)^{r}-n^{r}=r\xi^{r-1}$. This implies (see~\cite[Lemma 2.6]{lichen})
\begin{align}\label{lowerboundforgn}
|\Delta_f(n)|
\geq r\xi^{r-1}-1
=r\xi^{r}+\xi^{r}-1
=(r-1)\xi^{r-1}+\xi^{r-1}-1
>(r-1)n^{r-1}.
\end{align}

\noindent
Next, we express
\begin{align*}
\left|1+ue^{-n(\tau+iy)}\right|^2
=(1+ue^{-n\tau-niy})(1-ue^{-n\tau+niy})
=(1+ue^{n\tau})^2-2ue^{-n\tau}(1-\cos(ny)).
\end{align*}
\noindent
Using the above identities, we obtain
\begin{align*}
\left(\frac{|Q_{\sigma_r}(e^{\tau+iy},u)|}{Q_{\sigma_r}(e^{-\tau,u})}\right)^2
=&\prod_{n=1}^{\infty}\left(1-\frac{2ue^{-n\tau}(1-\cos(ny))}{(1+ue^{-n\tau})^2}\right)^{\Delta_f(n)}\\
\leq &\exp\left(\sum_{n=1}^{\infty}|\Delta_f(n)|\log\left|1-\frac{2ue^{-n\tau}(1-\cos(ny))}{(1+ue^{-n\tau})^2}\right|\right)\\
\leq&\exp\left(-\sum_{n=1}^{\infty}|\Delta_f(n)|\left|\frac{2ue^{-n\tau}(1-\cos(ny))}{(1+ue^{-n\tau})^2}\right|\right)\\
\leq&\exp\left(-\frac{2u}{(1+u)^2}\sum_{n=1}^{\infty}|\Delta_f(n)|e^{-n\tau}(1-\cos(ny))\right)\\
\leq&\exp\left(-\frac{2u(r-1)}{(1+u)^2}\sum_{n=1}^{\infty}n^{r-1}e^{-n\tau}(1-\cos(ny))\right).
%\leq&\exp\left(-\frac{2uk}{(1+u^2)}\sum_{n=1}^{\infty}n^{k+1-j}e^{-n\tau}(1-\cos(ny))\right).
\end{align*}

\noindent
Finally, applying Lemma~\ref{lichenlemma}, yields

\begin{align}\label{mainexpression}
\left(\frac{|Q_{\sigma_r}(e^{\tau+iy},u)|}{Q_{\sigma_r}(e^{-\tau,u})}\right)^2
\leq\exp\left(-\frac{2u(r-1)}{(1+u)^2}\varrho\left(\frac{e^{-\tau}}{(1-e^{-\tau})^{r}}-\frac{e^{-\tau}}{\left|1-e^{-\tau-iy}\right|^{r}}\right)\right).
\end{align}

%While the subsequent argument bears some resemblance to \cite[Lemma 7]{lipniketal}, we have considered the terms in the preceding expressions with more precision. 

\medskip

Beginning with a few standard trigonometric identities from~\cite{lichen}, we express
\begin{align}\label{trig1}
\left|1-e^{-\tau-i\theta}\right|^{r}
=&((1-e^{-\tau})^2+2e^{-\tau}(1-\cos\theta))^{\tfrac{r}{2}}\nonumber\\
\geq & (1-e^{-\tau})^{r}\left(1+\frac{2e^{-\tau}}{(1-e^{-\tau})^2}(1-\cos\theta_n)\right)^{\tfrac{r}{2}},
\end{align}
where $\theta_n\leq |\theta|\leq \pi$. By employing the Taylor series expansion of $\cos{\theta}$, we obtain
\begin{align*}
1-\cos\theta_n
=\frac{1}{2}\theta_n^2+O(\theta_n^4)
=\frac{1}{2}\tau^{2+{6r}/{7}}+O(\theta_n^{4+12r/7}),
\end{align*}
and similarly, $e^{-\tau}=1-\tau+O(\tau^2)$. Consequently, we derive
\begin{align*}
\frac{2e^{-\tau}}{(1-e^{-\tau})^2}(1-\cos\theta_n)
=\frac{2(1+O(\tau))(\tfrac{1}{2}\tau^{2+{6r}/{7}}+O(\theta_n^{4+12r/7}))}{\tau^2(1+O(\tau))}
=\tau^{{6r}/{7}}(1+O(\tau)).
\end{align*}

\medskip
\noindent
Substituting the above expression into~\eqref{trig1} gives
\begin{align*}
\left|1-e^{-\tau-i\theta}\right|^{r}
\geq (1-e^{-\tau})^{r}\left(1+\tau^{{6r}/{7}}(1+O(\tau))\right)^{\tfrac{r}{2}}\\
=(1-e^{-\tau})^{r}\left(1+\frac{r\tau^{{6r}/{7}
}}{2}+O\left(\tau^{{6r}/{7}+1}+\tau^{{12r}/{7}}\right)\right).
\end{align*}

\noindent
Therefore, the inner sum of~\eqref{mainexpression} simplifies to
\begin{align*}
\frac{e^{-\tau}}{(1-e^{-\tau})^{r}}-\frac{e^{-\tau}}{\left|1-e^{-\tau-iy}\right|^{r}}
&\geq\frac{e^{-\tau}}{(1-e^{-\tau})^{r}}\left(1-\frac{1}{1+\frac{r\tau^{{6r}/{7}}}{2}+O\left(\tau^{{6r}/{7}+1}+\tau^{{12r}/{7}}\right)}\right)\\
&=\frac{e^{-\tau}}{(1-e^{-\tau})^{r}}\left(\frac{r\tau^{{6r}/{7}}}{2}+O\left(\tau^{{6r}/{7}+1}+\tau^{{12r}/{7}}\right)\right)\\
&=\frac{1-O(\tau)}{\tau^{r}(1-O(\tau))}\left(\frac{r\tau^{{6r}/{7}}}{2}+O\left(\tau^{{6r}/{7}+1}+\tau^{{12r}/{7}}\right)\right)\\
&\geq\frac{r\tau^{-{r}/{7}}}{2}+O\left(\tau^{1-{r}/{7}}+\tau^{12r/7}\right)
\gg\frac{1}{4}\tau^{-\tfrac{r}{7}}.
\end{align*}

\noindent
Substituting this bound into~\eqref{mainexpression}, we conclude the proof.
\end{proof}

\medskip
\noindent
{\it Proof of Lemma~\ref{omega2sigmak}.} By the definition given in~\eqref{errortermintegral} for $f(n)=\sigma_r(n)$, the following inequity holds
\begin{align*}
|\Omega_{2,\sigma_r}|
\leq\frac{e^{n\tau}}{2\pi}\int_{\theta_n}^{\pi}\left|{Q}_{\sigma_r}(e^{-\tau-i\theta},u)\right|d\theta 
=\frac{e^{n\tau+F(\tau,u)}}{\pi}\int_{\theta_n}^{\pi}\frac{|Q_{\sigma_r}(e^{-\tau-i\theta},u)|}{Q_{\sigma_r}(e^{-\tau},u)}d\theta.
\end{align*}

\noindent
By employing Lemma~\ref{qlemmaforomega2sigmak}, we arrive at
\begin{align*}
|\Omega_{2,\sigma_r}|
\ll\exp\left(n\tau+F(\tau,u)-c_{r}\tau^{-r/7}\right),
\end{align*}
where $c_{r}$ is a constant that remains uniform in $u$, as noted.
\qed





\section{The Moment Generating Function}

In this section, we employ the estimates of $\Omega_{1,f}$ and $\Omega_{2,f}$ to derive
\begin{align*}
\mathcal{Q}_{f,n}(u)
=\frac{e^{n\tau+F(\tau,u)}}{\sqrt{2\pi B}}\left(1+O\left(\tau^{2r/7-1}\right)\right),
\end{align*}
where $r\geq 2$ is fixed and $B^2=F_{\gamma\gamma}(\tau,u)$. \\

We adopt the method outlined in~\cite[\S4.3]{lipniketal}, by considering the asymptotic expansion of the moment generating function associated with the random variable ${\varpi}_{f,n}$. Let $M_{f,n}(\theta)=\mathbb{E}(e^{(\varpi_{f,n}-\mu_{f,n})\theta/\nu_{f,n}})$ denote the moment generating function for real $\theta$, and let $\mu_{f,n}$ and $\nu_{f,n}$ represent the mean and standard deviation, as defined in~\eqref{meansigma} and~\eqref{variancesigma}, respectively. This leads to the following result.



\begin{lemma}\label{momentgeneratingfunctionlemma}
Let $r\geq 2$ be fixed and $f(n)=\sigma_r(n)$ be the generalised divisor function. Then as $n\to\infty$, for bounded $\theta$, the following holds
\begin{align*}
M_{f,n}(\theta)
=\exp\left(\frac{\theta^2}{2}+O\left(n^{-2r/7(r+2)+1/(r+2)}\right)\right).
\end{align*}
\end{lemma}

\begin{proof}
We begin by noting that
\begin{align}\label{asymtoticsofmomentgeneratingfunction}
M_{f,n}(\theta)
=&\exp\left(-\frac{\mu_{f,n}\theta}{\nu_{f,n}}\right)\frac{\mathcal{Q}_{f,n}(e^{\theta/\nu_{f,n}})}{\mathcal{Q}_{f,n}(1)}\nonumber\\
=&\sqrt{\frac{B^2(\tau(n,1),1)}{B^2(\tau(n,e^{\theta/\nu_{f,n}}),e^{\theta/\nu_{f,n}})}}\exp\Bigg(-\frac{\mu_{f,n}\theta}{\nu_{f,n}}+n\tau(n,e^{\theta/\nu_{f,n}})\nonumber\\
&\quad\quad\quad+{F}(\tau(n,e^{\theta/\nu_{f,n}}),e^{\theta/\nu_{f,n}})-n\tau(n,1)-{F}(\tau(n,1),1)+O(\tau^{2r/7-1})\Bigg).
\end{align}
\noindent
For simplicity of notation, we define $\tau=\tau(n,u)$ which is a function of $n$ and $u$ with respect to $u$. By employing the saddle-point solution given in~\eqref{saddlepointsolutionforsigmak}, we obtain (following the approach of~\cite{madwag})

\begin{align*}
\tau_u
=\tau_{u}(n,u)
=-\frac{{F}_{u\gamma}(\tau,u)}{{F}_{\gamma\gamma}(\tau,u)}
=\frac{\sum_{k=1}^{\infty}\frac{k\Delta_f(k)e^{k\tau}}{(e^{k\tau+u})^2}}{\sum_{k=1}^{\infty}\frac{uk^2\Delta_f(k)e^{kr}}{(e^{kr}+u)^2}},
\end{align*}
and similarly,
\begin{align*}
\tau_{uu}=\tau_{uu}(n,u)
=\frac{1}{{F}_{\gamma\gamma}(\tau,u)^3}\big(-{F}_{\gamma\gamma\gamma}(\tau,u){F}_{u\gamma}(\tau,u)^2
&+2{F}_{u\gamma\gamma}(\tau,u){F}_{u\gamma}(\tau,u){F}_{\gamma\gamma}(\tau,u)\\
&-{F}_{uu\gamma}(\tau,u){F}_{\gamma\gamma}(\tau,u)^2\big),
\end{align*}
and
\begin{align*}
\tau_{uuu}=\tau_{uuu}(n,u)
=&\frac{1}{{F}_{\gamma\gamma}(\tau,u)^5}\big(-{F}_{uuu\gamma}(\tau,u){F}_{\gamma\gamma}(\tau,u)^4\\
&+(3{F}_{uu\gamma\gamma}(\tau,u){F}_{u\gamma}(\tau,u)+3{F}_{uu\gamma}(\tau,u){F}_{u\gamma\gamma}(\tau,u)){F}_{\gamma\gamma}(\tau,u)^3\\
&+\big(-3{F}_{u\gamma\gamma\gamma}(\tau,u){F}_{u\gamma}(\tau,u)^2-6{F}_{u\gamma\gamma}(\tau,u)^2{F}_{u\gamma}(\tau,u)\\
&-3{F}_{uu\gamma}(\tau,u){F}_{\gamma\gamma\gamma}(\tau,u){F}_{u\gamma}(\tau,u)\big)F_{\gamma\gamma}(\tau,u)^2\\
&+\big({F}_{\gamma\gamma\gamma\gamma}(\tau,u){F}_{u\gamma}(\tau,u)^3+9{F}_{u\gamma\gamma}(\tau,u){F}_{\gamma\gamma\gamma}(\tau,u){F}_{u\gamma}(\tau,u)^2\big){F}_{\gamma\gamma}(\tau,u)\\
&-3{F}_{u\gamma}(\tau,u)^3{F}_{\gamma\gamma\gamma}(\tau,u)^2\big).
\end{align*}

\noindent
As it is evident from the preceding expression, it is necessary to estimate the partial derivative of ${F}_{\gamma}(\tau,u)$ with respect to $u$. To achieve this, we apply the estimate established in Lemma~\ref{maintermforsigmar} and proceed using a method analogous to that employed previously with the converse mapping. We provide a detailed illustration for ${F}_{u\gamma}(\tau,u)$; the remaining cases follow in a similar manner. For $j=1$, from~\eqref{fgammauderivativefinal} we write

\begin{align*}
{F}_{\gamma}(\gamma,u)
=\mathfrak{h}_{1,1}(\gamma,u)-\mathfrak{h}_{2,1}(\gamma,u)
=\sum_{n=1}^{\infty}\sigma_{r}(n+1)n\sum_{\ell=1}^{\infty}(-u)^{\ell}e^{-n\ell\gamma}-\sum_{n=1}^{\infty}\sigma_{r}(n)n\sum_{\ell=1}^{\infty}(-u)^{\ell}e^{-n\ell\gamma}.
\end{align*}

\noindent
Then the Mellin transform of $\mathfrak{h}_{1,1}(\gamma,u)-\mathfrak{h}_{2,1}(\gamma,u)$ is given by
\begin{align*}
&\mathcal{H}_{1,1}(s,u)-\mathcal{H}_{2,1}(s,u)
=O\Bigg(\zeta(r+1)\mathrm{Li}_{s}(-u)\Gamma(s)\bigg(\sum_{i=0}^{r}\binom{r}{i}\frac{\zeta(s-i-1)}{\zeta(s-i+r)}K_r(s-i-1)\\
&\quad\quad\quad\quad\quad\quad\quad\quad\quad+\sum_{i=0}^{r}\binom{r}{i}\frac{\zeta(\sigma-i-1)\zeta(\sigma-i-1+r)\zeta(\sigma-i+r)}{\zeta(2(\sigma-i-1+r))}\zeta(r)E_r(\sigma-i-1)\\
&\quad\quad\quad\quad\quad\quad\quad\quad\quad\quad\quad\quad\quad\quad\quad\quad\quad\quad\quad\quad\quad\quad\quad\quad\quad-\frac{\zeta(s-1)\zeta(s-1-r)}{\zeta(r+1)}\bigg)\Bigg),
\end{align*}
where $s=\sigma+it$ with $\sigma,t\in\mathbb{R}$, and $K_r(s)$ and $E_r(s)$ are defined in~\eqref{arsdefinition} and~\eqref{erdefinition}, respectively. Differentiating with respect to $u$, we obtain
\begin{align*}
&\frac{\partial}{\partial u}(\mathcal{H}_{1,1}(s,u)-\mathcal{H}_{2,1}(s,u))
=O\Bigg(\frac{1}{u}\mathrm{Li}_{s-1}(-u)\Gamma(s)\zeta(r+1)\bigg(\sum_{i=0}^{r}\binom{r}{i}\frac{\zeta(s-i-1)}{\zeta(s-i+r)}K_r(s-i-1)\\
&\quad\quad\quad\quad\quad\quad\quad\quad\quad\quad\quad\quad+\sum_{i=0}^{r}\binom{r}{i}\frac{\zeta(\sigma-i-1)\zeta(\sigma-i-1+r)\zeta(\sigma-i+r)}{\zeta(2(\sigma-i-1+r))}\zeta(r)E_r(\sigma-i-1)\\
&\quad\quad\quad\quad\quad\quad\quad\quad\quad\quad\quad\quad\quad\quad\quad\quad\quad\quad\quad\quad\quad\quad\quad\quad\quad\quad\quad-\frac{\zeta(s-1)\zeta(s-1-r)}{\zeta(r+1)}\bigg)\Bigg).
\end{align*}

\noindent
By applying Theorem~\ref{cmtheorem}, yields
\begin{align*}
F_{\gamma u}(\tau,u)
=& O\Bigg(\frac{1}{u}\zeta(r+1)\bigg(\sum_{i=0}^{r}\binom{r}{i}\mathrm{Li}_{i+1}(-u)\Gamma(i+2)\frac{K_r(1)}{\zeta(r+2)}\tau^{-(i+2)}\\
+&\sum_{i=0}^{r}\binom{r}{i}\mathrm{Li}_{i+1}(-u)\Gamma(i+2)\frac{\zeta(r+1)\zeta(r+2)\zeta(r)E_r(1)}{\zeta(2(r+1))}\tau^{-(i+2)}-\mathrm{Li}_{r+1}(-u)\Gamma(r+2)\tau^{-(r+2)}\bigg)\Bigg)\\
=& O\left(\frac{N(r)}{u}\mathrm{Li}_{r+1}(-u)\tau^{-(r+2)}\Gamma(r+2)\right)
=O(n),
\end{align*}
where $N(r)$ is defined in~\eqref{nrdefinition}. In a similar fashion, we derive the remaining partial derivatives and get
\begin{align*}
&F_{u}(\tau,u)
=\sum_{k=1}^{\infty}\frac{\Delta_f(k)}{e^{k\tau}+u}
=O\left(\frac{N(r)}{u}\mathrm{Li}_{r+1}(-u)\tau^{-(r+1)}\Gamma(r+1)\right)
=O(n^{\frac{r+1}{r+2}}),\\
&F_{uu}(\tau,u)
=-\sum_{k=1}^{\infty}\frac{\Delta_f(k)}{(e^{k\tau}+u)^2}
=O\left(\frac{N(r)}{u^2}(\mathrm{Li}_{r}(-u)-\mathrm{Li}_{r+1}(-u))\Gamma(r+1)\tau^{-(r+1)}\right)
=O(n^{\frac{r+1}{r+2}}),\\
&F_{\gamma\gamma}(\tau,u)
=\sum_{k=1}^{\infty}\frac{uk^2\Delta_f(k)e^{k\tau}}{(e^{k\tau}+u)^2}
=O\left(N(r)\mathrm{Li}_{r+2}(-u)\Gamma(r+3)\tau^{-(r+3)}\right)
=O(n^{\frac{r+3}{r+2}}),\\
&F_{uuu}(\tau,u)
=\sum_{k=1}^{\infty}\frac{2\Delta_f(k)}{(e^{k\tau}+u)^3}
=O(\tau^{-(r+1)})
=O(n^{\frac{r+1}{r+2}}),\\
&F_{uu\gamma}(\tau,u)
=\sum_{k=1}^{\infty}\frac{2k\Delta_f(k)e^{k\tau}}{(e^{k\tau}+u)^3}
=O(\tau^{-(r+2)})
=O(n),\\
&F_{u\gamma\gamma}(\tau,u)
=\sum_{k=1}^{\infty}\frac{k^2\Delta_f(k)e^{k\tau}(e^{k\tau}-u)}{(e^{k\tau}+u)^3}
=O(\tau^{-(r+3)})
=O(n^{\frac{r+3}{r+2}}),\\
&F_{\gamma\gamma\gamma}(\tau,u)
=-\sum_{k=1}^{\infty}\frac{uk^3\Delta_f(k)e^{k\tau}(e^{k\tau}-u)}{(e^{k\tau}+u)^3}
=O(\tau^{-(r+4)})
=O(n^{\frac{r+4}{r+2}}),\\
&F_{uuu\gamma}(\tau,u)
=-\sum_{k=1}^{\infty}\frac{6k\Delta_f(k)e^{k\tau}}{(e^{k\tau}+u)^4}
=O(\tau^{-(r+2)})
=O(n),\\
&F_{uu\gamma\gamma}(\tau,u)
=-\sum_{k=1}^{\infty}\frac{2k^2\Delta_f(k)e^{k\tau}(e^{k\tau}-u)}{(e^{k\tau}+u)^4}
=O(\tau^{-(r+3)})
=O(n^{\frac{r+3}{r+2}}),\\
&F_{u\gamma\gamma\gamma}(\tau,u)
=-\sum_{k=1}^{\infty}\frac{k^3\Delta_f(k)e^{k\tau}(e^{2k\tau}-4ue^{k\tau}+u^2)}{(e^{k\tau}+u)^4}
=O(\tau^{-(r+4)})
=O(n^{\frac{r+4}{r+2}}),\\
&F_{\gamma\gamma\gamma\gamma}(\tau,u)
=\sum_{k=1}^{\infty}\frac{uk^4\Delta_f(k)e^{k\tau}(e^{2k\tau}-4ue^{k\tau}+u^2)}{(e^{k\tau}+u)^4}
=O(\tau^{-(r+5)})
=O(n^{\frac{r+5}{r+2}}).
\end{align*}

\noindent
The above bounds imply that the derivatives $\tau_u$, $\tau_{uu}$, and $\tau_{uuu}$ are each $O(n^{-1/(r+2)})$, uniformly with respect to $u$. Expanding $\tau(n, e^{\theta/\nu_{f,n}})$ and $F(\tau(e^{\theta/\nu_{f,n}},n), e^{\theta/\nu_{f,n}})$ around $\theta=0$ yields
\begin{align*}
\tau(n,e^{\theta/\nu_{f,n}})-\tau(n,1)
=\tau_{u}(n,1)\frac{\theta}{\nu_{f,n}}+\frac{\tau_{u}(n,1)+\tau_{uu}(n,1)}{2}\left(\frac{\theta}{\nu_{f,n}}\right)^2+O\left(n^{-1/(r+2)}\frac{|\theta|^3}{\nu^3_{f,n}}\right)
\end{align*}
and
\begin{align*}
&F(\tau(e^{\theta/\nu_{f,n}},n), e^{\theta/\nu_{f,n}})-F((\tau,1),1)\\
=&F_{\gamma}(\tau(n,1),1)\tau_{u}(n,1)+F_{u}(\tau(n,1),1)\frac{\theta}{\nu_{f,n}}
+\frac{1}{2}\bigg(F_{\gamma}(\tau(n,1),1)(\tau_{u}(n,1)+\tau_{uu}(n,1))\\
&+F_{\gamma\gamma}(\tau(n,1),1)\tau_{u}(n,1)^2+2F_{u\gamma}(\tau(n,1),1)\tau_{u}(n,1)+F_{u}(\tau(n,1),1)\\
&+F_{uu}(\tau(n,1),1)\bigg)\left(\frac{\theta}{\nu_{f,n}}\right)^2+O\left(n^{\frac{r+1}{r+2}}\frac{|\theta|^3}{\nu^3_{f,n}}\right),
\end{align*}
respectively. Substituting these expansions into the exponential in~\eqref{asymtoticsofmomentgeneratingfunction}, we obtain
\begin{align*}
&\left(n\tau_u(n,1)+F_{\gamma}(\eta,1)\tau_u(n,1)+F_u(\eta,1)-\mu_{f,n}\right)\frac{\theta}{\nu_{f,n}}+\frac{1}{2}\bigg(n(\tau_u(n,1)+\tau_{uu}(n,1))+F_{\gamma}(\eta,1)(\tau_u(n,1)\\
&+\tau_{uu}(n,1))
+F_{\gamma\gamma}(\eta,1)\tau_u(n,1)^2+2F_{u\gamma}(\eta,1)\tau_u(n,1)+F_{u}(\eta,1)+F_{uu}(\eta,1)\bigg)\left(\frac{\theta}{\nu_{f,n}}\right)^2\\
&+O\left(n^{\frac{r+1}{r+2}}\frac{|\theta^3|}{\nu^3_{f,n}}+n^{-2r/7(r+2)+1/(r+2)}\right).
\end{align*}
Here, we have written $\eta = \tau(n, 1)$ for brevity. Recalling that $n = -F_{\gamma}(\eta, 1)$ and $\tau_u(n, 1) = -\dfrac{F_{u\gamma}(\eta, 1)}{F_{\gamma\gamma}(\eta, 1)}$, we may simplify the final expression to derive
\begin{align*}
&(F_{u}(\eta,1)-\mu_{f,n})\frac{\theta}{\nu_{f,n}}+\frac{1}{2}\bigg(F_{u}(\eta,1)+F_{uu}(\eta,1)-\frac{F_{u\gamma}(\eta,1)^2}{F_{\gamma\gamma}(\eta,1)}\bigg)\left(\frac{\theta}{\nu_{f,n}}\right)^2\\
&+O\left(n^{\frac{r+1}{r+2}}\frac{|\theta^3|}{\nu^3_{f,n}}+n^{-2r/7(r+2)+1/(r+2)}\right).
\end{align*}

\noindent
In a similar manner, we obtain that
\begin{align*}
\frac{B^2(\tau(n,1),1)}{B^2(\tau(n,e^{\theta/\nu_{f,n}}),e^{\theta/\nu_{f,n}})}
=1+O\left(\frac{|\theta|}{\nu_{f,n}}\right).
\end{align*}

\noindent
Thus, we derive the following asymptotic expression for the moment generating function in~\eqref{asymtoticsofmomentgeneratingfunction}:
\begin{align*}
M_{f,n}(\theta)
=&\exp\Bigg((F_{u}(\eta,1)-\mu_{f,n})\frac{\theta}{\nu_{f,n}}+\frac{1}{2}\bigg(F_{u}(\eta,1)+F_{uu}(\eta,1)-\frac{F_{u\gamma}(\eta,1)^2}{F_{\gamma\gamma}(\eta,1)}\bigg)\left(\frac{\theta}{\nu_{f,n}}\right)^2\\
&+O\left(n^{\frac{r+1}{r+2}}\frac{|\theta^3|}{\nu^3_{f,n}}+n^{-2r/7(r+2)+1/(r+2)}\right)\Bigg).
\end{align*}

\noindent
Recall that we choose $\mu_{f,n}$ and $\nu_{f,n}$ in~\eqref{meansigma} and~\eqref{variancesigma} such that
\begin{align}\label{meanvarianceformulatosolve}
&\mu_{f,n}
=F_{u}(\eta,1)
=\sum_{k=1}^{\infty}\frac{\Delta_f(k)}{e^{\eta k}+1}\quad\text{and}\quad\nonumber\\
&\nu^2_{f,n}
=F_{u}(\eta,1)+F_{uu}(\eta,1)-\frac{F_{u\gamma}(\eta,1)^2}{F_{\gamma\gamma}(\eta,1)}
=\sum_{k=1}^{\infty}\frac{\Delta_f(k)e^{\eta k}}{(e^{\eta k}+1)^2}-\frac{\left(\sum_{k=1}^{\infty}\frac{\Delta_f(k)ke^{\eta k}}{(e^{\eta k}+1)^2}\right)^2}{\sum_{k=1}^{\infty}\frac{\Delta_f(k)k^2e^{\eta k}}{(e^{\eta k}+1)^2}}.
\end{align}

\noindent
From the definitions of $\mu_{f,n}$ and $\nu^2_{f,n}$ in~\eqref{meansigma} and~\eqref{variancesigma} respectively, we deduce that
\begin{align*}
M_{f,n}(\theta)
=&\exp\left(\frac{\theta^2}{2}+O\left(n^{-(r+1)/(r+2)}|\theta|+n^{-(r+1)/(r+2)}|\theta|^3+n^{-2r/7(r+2)+1/(r+2)}\right)\right)\\
=&\exp\left(\frac{\theta^2}{2}+O\left(n^{-2r/7(r+2)+1/(r+2)}\right)\right),
\end{align*}
for bounded $\theta$.
\end{proof}


\subsection{Formulae for the mean and variance} 

By applying Lemma~\ref{momentgeneratingfunctionlemma} and Curtiss' theorem~\cite{curtiss}, we establish that the distribution of ${\varpi}_{f,n}$ is indeed asymptotically normal. In this section, we derive an upper bound for its mean and variance, providing an explicit formula for the constant as a function of $r$.

\medskip

From~\eqref{meanvarianceformulatosolve} we write
\begin{align*}
&\mu_{f,n}
=F_{u}(\eta,1)
=O\left(N(r)\mathrm{Li}_{r+1}(-1)\Gamma(r+1)/\eta^{(r+1)}\right)
\end{align*}
and 
\begin{align*}
\nu^2_{f,n}
=&F_{u}(\eta,1)+F_{uu}(\eta,1)-\frac{F_{u\gamma}(\eta,1)^2}{F_{\gamma\gamma}(\eta,1)}\\
=& O\left(\left(N(r)\mathrm{Li}_{r}(-1)\Gamma(r+1)-N(r)\frac{\mathrm{Li}^2_{r}(-1)\Gamma^2(r+2)}{\mathrm{Li}_{r+1}(-1)\Gamma(r+3)}\right)\frac{1}{\eta^{(r+1)}}\right).
\end{align*}

\noindent
Utilising the saddle-point solution in~\eqref{saddlepointsolutionforsigmak}, we write
\begin{align*}
\frac{1}{\eta}
=\frac{1}{\tau(n,1)}
=O\left(\frac{n^{\frac{1}{r+2}}}{{N(r)\mathrm{Li}_{r+2}(-1)\Gamma(r+2)}}\right)
\end{align*}
then 
\begin{align*}
\mu_{f,n}
=O\left(C_{\mu}(r)n^{\frac{r+1}{r+2}}\right)
\end{align*}
and
\begin{align*}
\nu^2_{f,n}
=O\left(C_{\nu}(r)n^{\frac{r+1}{r+2}}\right),
\end{align*}
where 
\begin{align}\label{meanconstant}
C_{\mu}(r)
\sim -N(r){\mathrm{Li}_{r+1}(-1)\Gamma(r+1)}%{(\mathrm{Li}_{r+2}(-1)\Gamma(r+2))^{r+1}}
=N(r)(1-2^{r+1})\Gamma(r+1)\zeta(r+1)
\end{align}
and 
\begin{align}\label{varianceconstant}
C_{\sigma}(r)
\sim& -N(r)\left({\mathrm{Li}_{r}(-1)\Gamma(r+1)-\frac{\mathrm{Li}^2_{r}(-1)\Gamma^2(r+2)}{\mathrm{Li}_{r+1}(-1)\Gamma(r+3)}}\right)\nonumber\\%{(\mathrm{Li}_{r+2}(-1)\Gamma(r+2))^{r+1}}\\
=&N(r)\zeta(r)(1-2^{-r})\left(\Gamma(r+1)-\frac{\Gamma(r+2)\zeta(r)}{\Gamma(r+3)\zeta(r+1)}\left(1-\frac{1}{2^{r+1}-1}\right)\right),
\end{align}
$N(r)$ is given in~\eqref{nrdefinition} and
\begin{align*}
-\mathrm{Li}_{r}(-1)
=(1-2^{-r})\zeta(r+1).
\end{align*}
\noindent
In this context, it is worth noting that the key constant of interest is $N(r)$, which emerges from a detailed analysis of the Dirichlet series with the gap function $\mathcal{D}_f(\Delta;s)$.

\subsubsection{The distribution of the tail}

We conclude the proof by bounding the tail of the distribution, following the approach of Lipnik et al.~\cite{lipniketal}, which in turn was adapted from Hwang~\cite{hwang}. For $\theta=o(n^{\frac{r+1}{6(r+2)}})$, we apply the Chernoff bound (see, for instance,~\cite{chernoff}) to obtain
\begin{align*}
\mathbb{P}\left(\frac{\varpi_{f,n}-\mu_{f,n}}{\sigma_{f,n}}\geq x\right)
\leq e^{-\theta x} M_{f,n}(\theta)
=e^{-t\theta+\theta^2/2}\left(1+O\left(n^{\frac{r+1}{2(r+2)}}(|\theta|+|\theta|^3)+n^{-2r/7(r+2)+1/(r+2)}\right)\right).
\end{align*}

\noindent
Setting $T=n^{\frac{r+1}{6(r+2)}}/\log n$ we have two cases. For the first case, let $x\leq T$, and we set $\theta=x$, yielding
\begin{align*}
\mathbb{P}\left(\frac{\varpi_{f,n}-\mu_{f,n}}{\sigma_{f,n}}\geq x\right)  
\leq e^{-x^2/2}\left(1+O((\log n)^{-3})\right).
\end{align*}

\noindent
Similarly, for $x\geq T$, we set $\theta=T$, obtaining 
\begin{align*}
\mathbb{P}\left(\frac{\varpi_{f,n}-\mu_{f,n}}{\sigma_{f,n}}\geq x\right)  
\leq e^{-Tx/2}\left(1+O((\log n)^{-3})\right).
\end{align*}

\noindent
A similar argument can be applied to estimate the probability $\mathbb{P}\left(\frac{\varpi_{f,n}-\mu_{f,n}}{\sigma_{f,n}}\leq -x\right)$.

\begin{comment}
\section{Proof of Theorem~\ref{clt}}\label{ramanujansumthmproof}

Theorem~\ref{clt} can be deduced from Theorem~\ref{cltforsigma} with slight modifications, as the Ramanujan sum $c_m(n)$ is considered with a fixed modulus $m$ satisfying conditions~(C.1), (C.2), and (C.3), which significantly simplifies the argument. Notably, these conditions ensure that $c_m(n)$ is a strictly non-decreasing function. Consequently, we can reformulate~\eqref{partitiondef} as

\begin{align*}
n=n_1c_m(n_1)+\cdots+n_rc_m(n_r),
\end{align*}
where $\{n_jc_m(n_j)\}_{j\in[1,r]}$ forms a strictly increasing sequence. Thus, we express 
\begin{align}\label{gnforcmn}
\Delta_f(n)=(n+1)c_m(n+1)-nc_m(n).    
\end{align}

\subsection{Saddle-point Solutions}
Following the same steps as in the preceding argument, we first estimate the integrands~\eqref{maintermintegral} and~\eqref{errortermintegral} for $f(n)=nc_m(n)$, where $m\geq 1$ is fixed and satisfies conditions~(C.1), (C.2), and (C.3).

\begin{lemma}\label{maintermlemmaforcqn}
Let $f(n)=nc_m(m)$ such that for fixed $m\geq 1$ satisfying conditions (C.1), (C.2), and (C.3). Define $B^2=F_{\gamma\gamma}$. Then, we have
\begin{align*}
\Omega_{f,1}
=\frac{e^{n\tau+F(\tau, u)}}{2\pi}\left(\int_{-\theta_n}^{\theta_n}\exp\left(-\frac{B^2}{2}\theta^2\right)d\theta\right)\left(1+O\left(\tau^{\frac{2}{7}}\right)\right),
\end{align*}
where
\begin{align*}
\int_{-\theta_n}^{\theta_n}\exp\left(-\frac{B^2}{2}\theta^2\right)d\theta 
=\frac{\sqrt{2\pi}}{B}+O\left(\tau^{-\frac{10}{7}}B^{-2}\exp\left(-\frac{\tau^{-1/7}}{2}\right)\right)
\end{align*}
uniformly in $u$ and the absolute constant depending on $m$.
\end{lemma}

\begin{proof}
Recalling~\eqref{saddlepointequationforsigmak} and~\eqref{gnforcmn} yields

\begin{align}\label{fgammau}
F(\gamma,u)
=&\sum_{n=1}^{\infty}\Delta_f(n)\log(1+ue^{-n\gamma})
=-\sum_{n=1}^{\infty}\Delta_f(n)\sum_{\ell=1}^{\infty}\frac{(-u)^{\ell}}{\ell}e^{-n\ell\gamma}\nonumber\\
=&-\sum_{n=1}^{\infty}(n(c_m(n+1)-c_m(n))+c_m(n+1))\sum_{\ell=1}^{\infty}\frac{(-u)^{\ell}}{\ell}e^{-n\ell\gamma}\nonumber\\
=&-\sum_{n=1}^{\infty}nc_m(n+1)\sum_{\ell=1}^{\infty}\frac{(-u)^{\ell}}{\ell}e^{-n\ell\gamma}
+\sum_{n=1}^{\infty}nc_m(n)\sum_{\ell=1}^{\infty}\frac{(-u)^{\ell}}{\ell}e^{-n\ell\gamma}\nonumber\\
&\quad\quad\quad\quad\quad\quad\quad\quad\quad\quad\quad\quad\quad\quad-\sum_{n=1}^{\infty}c_m(n+1)\sum_{\ell=1}^{\infty}\frac{(-u)^{\ell}}{\ell}e^{-n\ell\gamma}\nonumber\\
=&-F_{1}(\gamma,u)+F_{2}(\gamma,u)-F_{3}(\gamma,u).
\end{align}


Observe that addressing the second sum in~\eqref{fgammau} is somewhat more straightforward than tackling the first summand and the third summands. Therefore, we start with the analysis of the second summand. For $j\geq 0$, we derive

\begin{align}\label{f2derivative}
\frac{\partial^j }{\partial\gamma^j}F_2(\gamma,u) 
=(-1)^{j}\sum_{n=1}^{\infty}c_m(n)n^{j+1}\sum_{\ell=1}^{\infty}{(-u)^{\ell}}{\ell^{j-1}}e^{-n\ell\gamma}
\end{align}

Note that all the infinite sum in~\eqref{f2derivative} are of the form,

\begin{align*}
h_{2,j}(\gamma,u)
=\sum_{n=1}^{\infty}c_m(n)n^{j+1}\sum_{\ell=1}^{\infty}{(-u)^{\ell}}{\ell^{j-1}}e^{-n\ell\gamma}.
\end{align*}

Let $H_{2,j}(s,u)$ be the Mellin transform of $h_{2,j}(\gamma,u)$ with respect to $\gamma$. Accordingly, we express it as\footnote{Since the Dirichlet series involving the Ramanujan sum satisfies $\sum_{n\geq 1}\frac{c_m(n)}{n^s}=\zeta(s)\sum_{d|m}\mu(m/d)d^{1-s}$, it follows, by an argument analogous to that in Lemma~\ref{dirichletseriesidentity}, that $\sum_{d|m}\mu(m/d)d^{-k}=\frac{1}{m^k}\prod_{p|m}(1-p^k)$.}

\begin{align*}
H_{2,j}(s,u)
=&\sum_{n=1}^{\infty}c_m(n)n^{j+1}\sum_{\ell=1}^{\infty}(-u)^{\ell}\ell^{j-1}\int_{0}^{\infty}e^{-n\ell\gamma}\gamma^{s-1}d\gamma\\
=&\sum_{n=1}^{\infty}c_m(n)n^{j+1}\sum_{\ell=1}^{\infty}(-u)^{\ell}\ell^{j-1}\Gamma(s)(n\ell)^{-s}\\
=&\sum_{n=1}^{\infty}\frac{c_m(n)}{n^{s-j-1}}\sum_{\ell=1}^{\infty}\frac{(-u)^{\ell}}{\ell^{s-j+1}}\Gamma(s)
=\zeta(s-j-1)\left(\frac{1}{m^{s-j-2}}\prod_{p|m}(1-p^{s-j-2})\right)\mathrm{Li}_{s-j+1}(-u)\Gamma(s).
\end{align*}

Note that $H_{2,j}(s,u)$ converges for $\mathrm{Re}(s)>j+3$, with its sole pole in the interval $j+\frac{3}{2}\leq\mathrm{Re}(s)\leq j+3$ being the pole of $\zeta(s-j-1)$ at $s=j+2$. The Riemann zeta function and the polylogarithm exhibit polynomial growth, while the Gamma function decays exponentially along every vertical line in the complex plane. Consequently, we invoke Theorem~\ref{cmtheorem} to obtain
\begin{align*}
h_{2,j}(\gamma,u)
=\mathrm{Li}_{3}(-u)\Gamma(j+2)\gamma^{-(j+2)}+O\left(\gamma^{-j-\tfrac{3}{2}}\right).
\end{align*}

Observe that that $h_{2,j}(\gamma,u)$ is non-vanishing if and only if $m=1$. Substituting the above expression into~\eqref{f2derivative}, yields
\begin{align}\label{f2derivativefinalexpression}
\frac{\partial^j }{\partial\gamma^j}F_2(\gamma,u) 
=(-1)^{j}h_{2,j}(\gamma,u)+O(\gamma^{-j-2})
=(-1)^{j}\mathrm{Li}_{3}(-u)\Gamma(j+2)\gamma^{-(j+2)}+O(\gamma^{-j-1}).
\end{align}

We now turn our attention to the first summand in~\eqref{fgammau}. 
For $j\geq 0$, and applying Lemma~\ref{shiftedramanujansumformula}, we express
\begin{align}\label{f1ju}
\frac{\partial^j}{\partial\gamma^j}F_1(\gamma,u)
=&(-1)^{j}\sum_{n=1}^{\infty}c_{m}(n+1)n^{j+1}\sum_{\ell=1}^{\infty}{(-u)^{\ell}}{\ell^{j-1}}e^{-n\ell\gamma}\nonumber\\
=&(-1)^{j}\left(\frac{\mu(m)}{\varphi(m)}\sum_{n=1}^{\infty}c_m(n)n^{j+1}+\sum_{\substack{\chi(\bmod{m})\\ \chi\ne \chi_0}}\tau(\chi)\sum_{n=1}^{\infty}c_{\Bar{\chi}}(n)n^{j+1}\right)\sum_{\ell=1}^{\infty}{(-u)^{\ell}}{\ell^{j-1}}e^{-n\ell\gamma}\nonumber\\
=&(-1)^{j}\left(h_{1,j}(\gamma,u)+\Tilde{h}_{1,j}(\gamma,u)\right).
\end{align}

Applying the Mellin transform on $h_{1,j}(\gamma,u)$ with respect to $\gamma$, and following the argument of Lemma~\ref{dsklemma}, we obtain
\begin{align*}
H_{2,j}(s,u)
=&\frac{\mu(m)}{\varphi(m)}\sum_{n=1}^{\infty}c_m(n)n^{j+1}\sum_{\ell=1}^{\infty}{(-u)^{\ell}}{\ell^{j-1}}\Gamma(s)(n\ell)^{-s}\\
=&\frac{\mu^2(m)}{m^{s-j-1}}\zeta(s-j-1)\prod_{p|m}(p^{s-j-2}+p^{s-j-3}+\cdots+p)\mathrm{Li}_{s-j+1}(-u)\Gamma(s).
\end{align*}

Building on the previous arguments, we aim to invoke Theorem~\ref{cmtheorem}. Noting that $H_{2,j}(s,u)$ converges for $\mathrm{Re}(s)>j+3$ with its only pole in the interval $j+\frac{3}{2}\leq \mathrm{Re}(s)\leq j+3$ being the pole of $\zeta(s-j-1)$ at $s=j+2$. Therefore,
\begin{align}\label{hiforcmn}
h_{1,j}(\gamma,u)
=\frac{\mu^2(m)}{m}\gamma^{-(j+2)}\mathrm{Li}_{3}(-u)\Gamma(j+2)+O(\gamma^{-j-\frac{3}{2}}).
\end{align}

Similarly, we apply the Mellin transform on $\Tilde{h}_{1,j}(\gamma,u)$ with respect to $\gamma$, and following the argument of Lemma~\ref{gausssumlemma}, we obtain

...........


A similar argument to that in with
\begin{align*}
h_{3,j}(\gamma,u)
=\sum_{n=1}^{\infty}c_m(n+1)n^j\sum_{\ell=1}^{\infty}(-u)^{\ell}\ell^{j-1}e^{-n\ell\gamma},
\end{align*}

and compute its Mellin transform, given by
\begin{align*}
H_{3,j}(s,u)
=\frac{m}{\varphi(m)}L(s-j,\chi_0)\mathrm{Li}_{s-j+1}(-u)\Gamma(s),
\end{align*}

which has a pole at $s=j+1$. Following analogous steps, we deduce that
\begin{align*}
h_{3,j}(\gamma,u)
=\gamma^{-j+1}\prod_{p|m}\left(1-\frac{1}{p^{j+1}}\right)\mathrm{Li}_{2}(-u)\Gamma(j+1)+O(\gamma^{-j-1}).
\end{align*}

This yields
\begin{align}\label{f3jufinal}
\frac{\partial^j}{\partial\gamma^j}F_3(\gamma,u)
=(-1)^j\mathrm{Li}_{2}(-u)\Gamma(j+1)\gamma^{-(j+1)}\prod_{p|m}\left(1-\frac{1}{p^{j+1}}\right)+O(\gamma^{-j-1}).
\end{align}


Therefore, by the linearity property of derivatives and~\eqref{fgammau}, we can express
\begin{align*}
\frac{\partial^j}{\partial\gamma^j}F(\gamma,u)
=-\frac{\partial^j}{\partial\gamma^j}F_1(\gamma,u)+\frac{\partial^j}{\partial\gamma^j}F_2(\gamma,u)-\frac{\partial^j}{\partial\gamma^j}F_3(\gamma,u).
\end{align*}

Note that~\eqref{f2derivativefinalexpression} and~\eqref{f1jufinal} will cancel each other. Therefore, only~\eqref{f3jufinal} will contribute to the main term, and hence we write

\begin{align}\label{finalderivativexpression}
\frac{\partial^j}{\partial\gamma^j}F(\gamma,u)
=\frac{\partial^j}{\partial\gamma^j}F_3(\gamma,u)
=(-1)^{j+1}\mathrm{Li}_{2}(-u)\Gamma(j+1)\gamma^{-(j+1)}\prod_{p|m}\left(1-\frac{1}{p^{j+1}}\right)+O(\gamma^{-j-1}).
\end{align}


Thus for the saddle-point $n$, we write
\begin{align}\label{saddlepointsoltuion}
n=-F_{\gamma}(\tau,u)
=-\mathrm{Li}_{2}(-u)\Gamma(2)\tau^{-2}\prod_{p|m}\left(1-\frac{1}{p^{2}}\right)+O(\tau^{-2}).
\end{align}

Now, for the second derivative we express
\begin{align*}
B^2
=F_{\gamma\gamma}(\tau,u)
=-\mathrm{Li}_{2}(-u)\Gamma(3)\tau^{-3}\prod_{p|m}\left(1-\frac{1}{p^{3}}\right)+O(\tau^{-3}),
\end{align*}

whereas for the third derivative we write
\begin{align}\label{thirdderivative}
\frac{\partial^3}{\partial\theta^3}F(\tau+i\theta,u)
=-i\sum_{n=1}^{\infty}\frac{n^3\Delta_f(n)u^{-1}e^{n(\tau+i\theta)}(1-u^{-1}e^{n(\tau+i\theta)})}{(u^{-1}e^{n(\tau+i\theta)}+1)^3}.
\end{align}

We adopt the argument outlined in~\cite[Lemma 4]{lipniketal}. Let $n_0=\tau^{-(1+c)}$ for some constant $c>0$, and for simplicity, we set $v=u^{-1}$. The summand is then divided into two parts based on whether $n\leq n_0$ or $n>n_0$. For the sum over $n>n_0$, we obtain

\begin{align}\label{thirdderivativesimplification}
&\left|-i\sum_{n>n_0}\frac{n^3\Delta_f(n)ve^{n(\tau+i\theta)}(1-ve^{n(\tau+i\theta)})}{(ve^{n(\tau+i\theta)}+1)^3}\right|
\leq\sum_{n>n_0}\frac{n^3\Delta_f(n)ve^{n\tau}(1+ve^{n\tau})}{|ve^{n(\tau+i\theta)}+1|^3}\nonumber\\
\leq&\sum_{n>n_0}\frac{n^3\Delta_f(n)ve^{n\tau}(1+ve^{n\tau})}{(ve^{n\tau}-1)^3}
=O\left(\sum_{n>n_0}\frac{n^3\Delta_f(n)}{e^{n\tau}}\right)
=O\left(\tau^{-4}\right).
\end{align}

For the final case, we first observe that
\begin{align*}
\left|1+ve^{n(\tau+i\theta)}\right|  
\geq \left(1+ve^{n\tau}\right)\cos\left(\frac{n\theta}{2}\right).
\end{align*}

Utilising the aforementioned identity, we express
\begin{align*}
&\left|-i\sum_{n\leq n_0}\frac{n^3\Delta_f(n)ve^{n(\tau+i\theta)}(1-ve^{n(\tau+i\theta)})}{(ve^{n(\tau+i\theta)}+1)^3}\right|
\leq\sum_{n\leq n_0}\frac{n^3\Delta_f(n)ve^{n\tau}(1+e^{n\tau})}{|1+ve^{n(\tau+i\theta)}|^3}\\
\leq&\sum_{n\leq n_0}\frac{n^3\Delta_f(n)ve^{n\tau}(1+e^{n\tau})}{(1+ve^{n\tau})^3}\left(1+O((n\theta)^2)\right)
\leq\sum_{n\leq n_0}\frac{n^3\Delta_f(n)}{ve^{n\tau}}\left(1+O((n\theta)^2)\right)\\
\leq&\sum_{n\geq 1}\frac{n^3\Delta_f(n)}{ve^{n\tau}}+O\left(\sum_{n\geq 1}\frac{n^5\Delta_f(n)\theta^2}{e^{n\tau}}\right).
\end{align*}

Employing the Mellin transform and the converse mapping (refer, for instance, to Theorem~\ref{cmtheorem}), we arrive at
\begin{align*}
\left|-i\sum_{n\leq n_0}\frac{n^3\Delta_f(n)u^{-1}e^{n(\tau+i\theta)}(1-u^{-1}e^{n(\tau+i\theta)})}{(u^{-1}e^{n(\tau+i\theta)}+1)^3}\right|
=O\left(\tau^{-4}+\tau^{-6}\theta^2\right).
\end{align*}

By combining this with~\eqref{thirdderivativesimplification}, we deduce that for $|\theta|\leq \tau^{13/14}$, the following holds.
\begin{align*}
\frac{\partial^3}{\partial\theta^3}F(\tau+i\theta,u)
=O\left(\tau^{-4}\right)
\end{align*}



In conclusion, this leads to the expansion of~\eqref{saddlepointequationforcqn} with
\begin{align*}
F(\tau+i\theta,u)
=F(\tau,u)-in\theta-\frac{B^2}{2}\theta^2+O(\tau^{-4}\theta^3).
\end{align*}

Finally, for the integral $\Omega_{1,c_m}$, we derive the following

\begin{align}\label{saddlepointintegral}
\frac{e^{n\tau}}{2\pi}\int_{-\theta_n}^{\theta_n}\exp(in\theta+F(\tau+i\theta,u))d\theta
=&\frac{e^{n\tau+F(\tau,u)}}{2\pi}\int_{-\theta_n}^{\theta_n}\exp\left(-\frac{B^2}{2}\theta^2+O\left(\tau^{2/7}\right)\right)d\theta\nonumber\\
=&\frac{e^{n\tau+F(\tau,u)}}{2\pi}\left(\int_{-\theta_n}^{\theta_n}\exp\left(-\frac{B^2}{2}\theta^2\right)d\theta\right)+\left(1+O\left(\tau^{2/7}\right)\right).
\end{align}

As the last step, we transform the integral into a Gaussian integral, which results in

\begin{align}\label{gaussianintegral}
\int_{-\theta_n}^{\theta_n}\exp\left(-\frac{B^2}{2}\theta^2\right)d\theta
=&\int_{-\infty}^{\infty}\exp\left(-\frac{B^2}{2}\theta^2\right)d\theta-2\int_{\theta_n}^{\infty}\exp\left(-\frac{B^2}{2}\theta^2\right)d\theta\nonumber\\
=&\frac{\sqrt{2\pi}}{B}+O\left(\int_{\theta_n}^{\infty}\exp\left(-\frac{B^2\tau^{1+3/7}}{2}\theta^2\right)d\theta\right)\nonumber\\
=&\frac{\sqrt{2\pi}}{B}+O\left(\tau^{-10/7}B^{-2}\exp\left(-\frac{-\tau^{-1/7}}{2}\right)\right). 
\end{align}



\end{proof}







 

\begin{lemma}\label{omega2cmn}
For fixed $m\geq 1$ sufficiently large $n$ with $f=c_m(n)$, the error term is bounded as follows
\begin{align*}
|\Omega_{2,c_m}|
=O\left(\exp\left(n\tau+F(\tau,u)-\mathfrak{c}_{m}\tau^{-k/7}\right)\right),    
\end{align*} 
where $\mathfrak{c}_{m}$ is a constant uniformly in $u$.
\end{lemma}

In order to prove Lemma~\ref{omega2cmn}, we need to derive an analogous statement to Lemma~\ref{qlemmaforomega2sigmak}.

\begin{lemma}\label{qlemmaforomega2cmn}
For fixed $m\geq 1$ and any real $y$ satisfying $\theta_0\leq |y|\leq \pi$, we have
\begin{align*}
\frac{|Q_{\sigma_k}(e^{\tau+iy},u)|}{Q_{\sigma_k}(e^{-\tau,u})}
\leq\left(\exp\left(-\frac{2uk}{(1+u)^2}\varrho\left(\frac{1}{4}\tau^{-\frac{k+1}{7}}\right)\right)\right),
\end{align*}
where $\varrho$ is a constant depends on $k$.    
\end{lemma}

\begin{proof}
    
\end{proof}


\subsection{Moment Generating Function}

\subsubsection{Formulae for the Mean and Variance}









\appendix

\section{Arithmetical identities}



\begin{lemma}\label{lfunctionavaragebound}
Let 
\begin{align*}
\Omega_{1}(s)
=\esum_{\substack{\chi(\bmod{m})\\ \chi\ne\chi_0}}L(s,\chi), \quad\text{and}\quad
\Omega_{2}(s,k)
=\osum_{\substack{\chi(\bmod{m})\\ \chi\ne\chi_0}}L(s,\chi)
\end{align*}
where $\chi$ is a primitive Dirichlet character and $\mathrm{Re}(s)>1$. Then for integer $k\geq 1$,
\begin{align}
|\Omega_{1}(2k)|
\ll\zeta(2k)\varphi(m)\sqrt{m}, \quad\text{and}\quad
|\Omega_{2}(2k+1)|
\ll \zeta(2k+1)\varphi(m)\sqrt{m}.
\end{align}
\end{lemma}

\begin{proof}
Recall that for an even primitive Dirichlet character $\chi(\bmod{m})$ and for $s=2k$ with an integer $k\geq 1$, the associated Dirichlet $L$-function satisfies 

\begin{align*}
L(2k,\chi) 
=\frac{(-1)^{k-1}2^{2k-1}\pi^{2k}\tau(\chi)}{m(2k)!}\sum_{a=1}^{m}\bar{\chi}(a)B_{2k}\left(\frac{a}{m}\right),
\end{align*}
where $B_k(x)$ denotes the Bernoulli polynomial (see, for example~\cite[Problems 14-15]{montgomeryvaughanbook}). Note that $L(2k,\chi)$ is entire for $k\geq 1$ since $\chi$ is primitive. Moreover, by employing the Fourier expansion
\begin{align*}
B_{2k}\left(\frac{a}{m}\right)
=\frac{(2k)!}{(2\pi i)^{2k}}\sum_{n=1}^{\infty}\frac{e^{2\pi ian/m}}{n^{2k}},
\end{align*}

we deduce that
\begin{align}\label{someidentity}
\left|\sum_{a=1}^{m}\bar{\chi}(a)B_{2k}\left(\frac{a}{m}\right)\right|
\ll 2m\zeta(2k)\frac{(2k)!}{(2\pi)^{2k}}.
\end{align}

Applying~\eqref{someidentity}, and recalling $|\tau(\chi)|=\sqrt{m}$ yields

\begin{align*}
|L(2k,\chi)|
=&\left|\frac{(-1)^{k-1}2^{2k-1}\pi^{2k}\tau(\chi)}{m(2k)!}\sum_{a=1}^{m}\bar{\chi}(a)B_{2k}\left(\frac{a}{m}\right)\right|\\
\ll& \left|\frac{(-1)^{k-1}2^{2k-1}\pi^{2k}\tau(\chi)}{m(2k)!} \cdot 2\zeta(2k)m\frac{(2k)!}{(2\pi)^{2k}}\right|
\ll \sqrt{m}\zeta(2k).
\end{align*}

Finally, by applying Cauchy's inequality together with
\begin{align*}
%\esum_{\substack{\chi(\bmod{m})\\\chi\ne\chi_0}}1
|\{\chi(\bmod{m}): \chi\ne \chi_0, \chi(-1)=1\}|=\frac{\varphi(m)}{2},    
\end{align*}
at $s=2k$, we obtain
\begin{align*}
\left|\quad\esum_{\substack{\chi(\bmod{m})\\ \chi\ne\chi_0}}L(2k,\chi)\right|
\leq \esum_{\substack{\chi(\bmod{m})\\ \chi\ne\chi_0}} |L(2k,\chi)|
\ll\frac{1}{2}\varphi(m)\sqrt{m}\zeta(2k).
\end{align*}

Similarly, for an odd primitive Dirichlet character $\chi(\bmod{m})$ satisfying $\chi(-1)=-1$ and for $s=2k+1$ with integer $k\geq 1$, we employ the identity
\begin{align*}
L(2k+1,\chi)
=\frac{i(-1)^k2^{2k}\pi^{2k+1}\tau(\chi)}{(2k+1)!m}\sum_{a=1}^{m}\bar{\chi}(a)B_{2k+1}\left(\frac{a}{q}\right),
\end{align*}
which allows us to derive an upper bound for $\Omega_2(s)$ with an analogous argument, thereby concluding the lemma.
\end{proof}

\begin{lemma}\label{mobiussquare}
For any positive integer $m\geq 1$, we have
\begin{align*}
\mu(m)(-1)^{\omega(m)}=\mu^2(m).    
\end{align*}
\end{lemma}

\begin{proof}
Let $m=p_1^{\alpha_1}p_2^{\alpha_2}\cdots p_r^{\alpha_r}$ denote the canonical factorization of $m$. Then, the L.H.S of the given identity can be written as
\begin{align}\label{mobiussquareexplained}
\mu(m)(-1)^{\omega(m)}
\begin{cases}
1 & \text{if $m=1$,}\\
1 & \text{if $m=\prod_{i=1}^{j}p_i$ for $1\leq j\leq r$,}\\
0 & \text{otherwise}.
\end{cases}
\end{align}

Note that all the terms involving higher order primes vanish due to the M\"{o}bius function and~\eqref{mobiussquareexplained} matches with the definition of $\mu^2(m)$, thus concluding the proof. 
\end{proof}
\end{comment}

\medskip

\medskip
\noindent
{\it{Acknowledgement}.} We are grateful to Kyle Pratt and Alexandru Zaharescu for their valuable suggestions and many insightful discussions. A special thanks to Oleksiy Klurman for his comments on an earlier version of this paper. The first author is supported by EPSRC doctoral fellowship \#EP/T518049/1.

\medskip
\noindent
{\it{Rights Retention}}. For the purpose of open access, the author has applied a Creative Commons Attribution (CC BY) licence to any Author Accepted Manuscript version arising from this submission.



\begin{thebibliography}{11}

\bibitem{alexandru}
E.~Alkan and A.~H. Ledoan and A.~Zaharescu.
\newblock {Asymptotic behavior of the irrational factor}.
\newblock \textit{Acta Mathematica Hungarica}, 13(3):293--305, 2008.

\bibitem{apostol}
T.~M. Apostol.
\newblock \textit{Introduction to Analytic Number Theory}.
\newblock Springer New York, NY, 1 edition, 2013.

\bibitem{baluluca}
R.~Balasubramanian and F.~Luca.
\newblock {On the number of factorizations of an integer}.
\newblock \textit{Integers}, 11(2):139--143, 2011.

\bibitem{mobius}
D.~Basak and N.~Robles and A.~Zaharescu.
\newblock {Exponential sums over M{\"{o}}bius convolutions with applications to
  partitions}.
\newblock \textit{To appear in Canadian Journal of Mathematics}, 2025.

\bibitem{gregmartin}
M.~A. Bennett and G.~Martin and K.~O'Bryant and A.~Rechnttzer.
\newblock {Counting zeros of dirichlet $L$-functions}.
\newblock \textit{Mathematics of Computations}, 80(329):1455--1482, 2021.


\bibitem{chernoff}
H.~Chernoff.
\newblock {A measure of asymptotic efficiency for tests of a hypothesis based on
  the sum of observations}.
\newblock \textit{Ann. Math. Statistics}, 23:493--507, 1952.

\bibitem{curtiss}
J.~H. Curtiss.
\newblock {A note on the theory of moment generating functions}.
\newblock \textit{Ann. Math. Statistics}, 13:430--433, 1942.

\bibitem{taylor}
T.~Daniels.
\newblock {Bounds on the M{\"{o}}bius-signed partition numbers}.
\newblock \textit{Ramanujan J.}, 65:81--123, 2024.

\bibitem{taylor2}
T.~Daniels.
\newblock {Legendre-signed partition numbers}.
\newblock \textit{J. Math. Anal. Appl.}, 542(1):128717, 2025.

\bibitem{madhu}
M.~Das.
\newblock {Exponential sums with additive coefficients and its consequences to
  weighted partitions}.
\newblock \textit{Canadian Math J.}, vol.~To appear, p.~41, 2025.

\bibitem{semiprime}
M.~Das and N.~Robles and A.~Zaharescu and D.~Zeindler.
\newblock {Partitions into semiprimes}.
\newblock \textit{Pre-prints: \url{https://arxiv.org/abs/2212.12489}}, p.~37,
  2022.

\bibitem{debten}
G.~Debruyne and G.~Tenenbaum.
\newblock {The saddle-point method for general partition functions}.
\newblock \textit{Indag. Math.}, 31:728--738, 2020.

\bibitem{erdoslehner}
P.~Erd{\H{o}}s and J.~Lehner.
\newblock {The distribution of the number of summands in the partitions of a
  positive integer}.
\newblock \textit{Duke Math. J.}, 8:335--345, 1941.

\bibitem{flajoletetal}
P.~Flajolet and X.~Gourdon and P.~Dumas.
\newblock {Mellin transforms and asymptotics: harmonic sums}.
\newblock \textit{Special volume on mathematical analysis of algorithms},
  pp.~3--58, 1995.

\bibitem{hardyramanujan}
G.~H. Hardy and S.~Ramanujan.
\newblock {Asymptotic formulae in combinatory analysis}.
\newblock \textit{Proceedings of the London Mathematical Society.}, 17:75--115,
  1918.

\bibitem{hwang}
H.~K. Hwang.
\newblock {Limit theorems for the number of summands in integer partitions}.
\newblock \textit{J. Combin. Theory Ser. A}, 96(1):89--126, 2001.

\bibitem{iwanieckowlaski}
H.~Iwaniec and E.~Kowalski.
\newblock \textit{{Analytic Number Theory}}, volume {American Mathematical
  Society colloquium publications}.
\newblock American Mathematical Society, 2004.
\newblock ISBN 9780821836330.

\bibitem{klurman}
O.~Klurman.
\newblock {Correlations of multiplicative functions and applications}.
\newblock \textit{Compositio Math.}, 153(8):1622--1657, 2017.

\bibitem{kratzel}
E.~Kr{\"{a}}tzel.
\newblock {Zahlentheorie}.
\newblock \textit{Berlin: VEB Deutscher Verlag der Wissenschaften}, p.~130, 1981.


\bibitem{lichen2}
Y.~L. Li and Y.~G. Chen.
\newblock {On the {$r$}-th root partition function}.
\newblock \textit{Taiwanese J. Math.}, 20(3):545--551, 2016.

\bibitem{lichen}
Y.~L. Li and Y.~G. Chen.
\newblock {On the {$r$}-th root partition function, II}.
\newblock \textit{J. Number Theory}, 188:392--409, 2018.

\bibitem{lipniketal}
G.~Lipnik and M.~Madritsch and R.~Tichy.
\newblock {A central limit theorem for integer partitions into small powers}.
\newblock \textit{Monatsh Math J.}, 203:149--173, 2024.

\bibitem{madwag}
M.~Madritsch and S.~Wagner.
\newblock {A central limit theorem for integer partitions}.
\newblock \textit{Monatsh. Math.}, 161(1):85--114, 2010.

\bibitem{sacha}
A.~Mangerel.
\newblock {Gap problems for integer-valued multiplicative functions}.
\newblock \textit{Pre-prints: arxiv {\url{https://arxiv.org/abs/2311.11636}}},
  pp.~1--45, 2023.

\bibitem{montgomeryvaughan}
H.L. Montgomery and R.C. Vaughan.
\newblock {Exponential Sums with Multiplicative Coefficients}.
\newblock \textit{Inventiones mathematicae}, 43:69--82, 1977.

\bibitem{montgomeryvaughanbook}
H.~L. Montgomery and R.~C. Vaughan.
\newblock \textit{Multiplicative Number Theory I}.
\newblock Cambridge University Press, 2006.

\bibitem{ramare}
O.~Ramar{\'{e}}.
\newblock {Accurate computations of Euler products over primes in arithmetic
  progressions}.
\newblock \textit{Funct. Approx. Comment. Math.}, 65(1):33--45, 2021.

\bibitem{ramanujan}
S.~Ramanujan.
\newblock {On certain arithmetical functions}.
\newblock \textit{Transactions of the Cambridge Philosophical Society},
  XXII(9):159--184, 1916.

\bibitem{richmond1}
L.~B. Richmond.
\newblock {The moments of partitions, I}.
\newblock \textit{Acta arithmetica}, 26(XXVI):411--425, 1975.

\bibitem{richmond2}
L.~B. Richmond.
\newblock {The moments of partitions, II}.
\newblock \textit{Acta arithmetica}, 28(XXVIII):229--243, 1975.


\bibitem{titchmarsh}
E.C. Titchmarsh.
\newblock \textit{The theory of functions}.
\newblock Oxford University Press, London, 2 edition, 1939.

\bibitem{vaughan}
R.~C. Vaughan.
\newblock {On the number of partitions into primes}.
\newblock \textit{Ramanujan J}, 15:109--121, 2008.

\end{thebibliography}

\end{document}
